\documentclass[12pt,a4paper]{article}

\usepackage{viu-mrob-midreport}
\usepackage[authoryear,round,sort]{natbib}
\usepackage{algorithm}
\usepackage{algpseudocode}
\usepackage{multicol}
\usetikzlibrary{shapes.geometric,positioning,patterns}
\usetikzlibrary{arrows.meta}

%--- Metadata
\setviutitle{Coordinación distribuida local de múltiples AGV para el transporte
  \mbox{cooperativo} de cargas heterogéneas}
\setviustudent{Jorge Luis Mayorga Taborda}
\setviututor{José Ignacio Iñíguez Amigot}
\setviuedition{2025--26}
\setviudate{mayo de 2026}
\setviuheadertitle{J.~L.~Mayorga $\cdot$ Coordinación distribuida de coaliciones multi-AGV}

%--- Comandos matemáticos
\newcommand{\RR}{\mathbb{R}}
\newcommand{\NN}{\mathbb{N}}
\newcommand{\simplex}{\Delta}
\newcommand{\pos}[1]{\left[#1\right]_{+}}
\newcommand{\norm}[1]{\left\|#1\right\|}
\newcommand{\xhat}{\hat{x}}
%--- Operadores. Se usan \min y \max estándar (sin acentos en subíndices) para
%    evitar artefactos de renderizado; en notación matemática min/max son estándar.
\DeclareMathOperator*{\argmax}{arg\,max}

\begin{document}

%--- Portada (pag 1, sin numero visible; titlepage env resetea contador)
\hypersetup{pageanchor=false}
\makeviucover
\hypersetup{pageanchor=true}

%--- Preliminares en numeración romana (resumen, índices, nomenclatura);
%    la portada no se numera. El cuerpo pasa a arábigo desde la Introducción.
\pagenumbering{roman}

%--- Pagina en blanco (pag 2, con header pero sin contenido)
\thispagestyle{frontmatter}
\null
\clearpage

%--- Resumen (pag 3)
\begin{viuabstract}
El transporte cooperativo de cargas heterogéneas en almacenes automatizados plantea
decisiones acopladas: qué robots atienden cada carga, cuántos colaboran y cómo se
desplazan hacia recogida y entrega. Ninguna la decide un coordinador central. La
pregunta de este informe es si una flota de AGV puede resolverlas con solo comunicación
local, estimando la ocupación de las tareas, asignándose a ellas y formando coaliciones
de cardinalidad mínima en un mismo bucle distribuido.

La validación se construye por niveles. Una primera línea base usa consenso dinámico de
media (DAC) con una regla voraz: sirve para ver si estimar bien la ocupación basta para
formar coaliciones. Sobre ella se contrastan dos dinámicas poblacionales, el replicador
y Smith, bajo idéntica estimación. La propuesta adopta Smith como política de revisión,
una sigmoide decreciente para el requisito mínimo de cardinalidad y una tasa de revisión
espacialmente modulada que frena las reasignaciones tardías.

El análisis formal cubre solo el caso homogéneo y balanceado con $\Psi \equiv 1$. Ahí el
payoff induce una estructura de juego de congestión y se prueba la estabilidad local del
equilibrio objetivo. Al añadir el descuento espacial $\Psi(d)$, la comunicación local y
la navegación reactiva, esa garantía deja de aplicar directamente; el sistema pasa a
tratarse como una versión perturbada y se mide en simulación. La campaña cuantitativa se
ejecutará en Python con seis escenarios y al menos treinta repeticiones por combinación.
CoppeliaSim queda para comprobar la plausibilidad cinemática con robots Pioneer~3-DX. La
manipulación cooperativa por contacto queda fuera del alcance.

\keywords{coordinación distribuida, formación de coaliciones, consenso dinámico,
  dinámicas poblacionales, transporte cooperativo multi-AGV.}

\vspace{0.4em}
\noindent{\bfseries\color{VIUOrange}Keywords:}\enspace distributed coordination,
coalition formation, dynamic average consensus, population dynamics, cooperative
multi-AGV transport.\par
\end{viuabstract}

%--- Indices y nomenclatura (siguen numeracion arabica)
\clearpage
\tableofcontents

\clearpage
\listoffigures

\clearpage
\listoftables

%% Nomenclatura
\clearpage
\section*{Nomenclatura}
\addcontentsline{toc}{section}{Nomenclatura}
\label{sec:nomenclatura}

{\small Símbolos agrupados en cinco bloques (índices, variables de estado,
parámetros, funciones y métricas). Entre paréntesis o corchetes se indica el
dominio o la unidad cuando procede.}

\begingroup
\scriptsize
\setlength{\parskip}{0pt}
\setlength{\tabcolsep}{3pt}
\setlength{\extrarowheight}{0pt}
\setlength{\aboverulesep}{1pt}
\setlength{\belowrulesep}{1pt}
\renewcommand{\arraystretch}{1.0}
\setlength{\columnsep}{1.2em}

% Cada bloque: etiqueta naranja + tabla de 2 columnas (simbolo | descripcion),
% envuelto en minipage para que no se parta entre columnas.
\newcommand{\nomblock}[2]{%
  \noindent\begin{minipage}{\linewidth}
  {\bfseries\color{VIUOrange}#1}\par\nobreak\smallskip
  \begin{tabularx}{\linewidth}{@{}l@{\hskip 6pt}X@{}}
  \toprule
  #2
  \bottomrule
  \end{tabularx}
  \end{minipage}\par\medskip}

\begin{multicols}{2}

\nomblock{Índices y conjuntos}{%
$i$               & Índice de robot ($\{1,\ldots,N\}$) \\
$k$               & Índice de tarea; $k{=}0$ inactividad ($\{0,\ldots,K\}$) \\
$t$               & Paso temporal discreto ($\mathbb{N}$) \\
$\mathcal{S}$     & Conjunto de estrategias ($\{0,\ldots,K\}$) \\
$\mathcal{N}_i(t)$& Vecinos de $i$ en el grafo de comunicación \\
$\mathcal{C}_k(t)$& Coalición asignada a la tarea $k$ \\
}

\nomblock{Variables de estado}{%
$\mathbf{q}_i$        & Pose $(x,y,\theta)$ [m, m, rad] \\
$\mathbf{q}_i^{xy}$   & Proyección plana de la pose [m] \\
$p_{ik}$              & Probabilidad mixta de $i$ sobre $k$ ($[0,1]$) \\
$\hat{x}_{ik}$        & Fracción estimada de robots en $k$ ($[0,1]$) \\
$\tilde{x}_k$         & Ocupación empírica pura $a_k/N$ ($[0,1]$) \\
$a_k$                 & Ocupación asignada pura $\sum_i\mathbf{1}_{[s_i=k]}$ \\
$\mathrm{count}_k$    & Asignados y presentes en $r_{\mathrm{det}}$ ($=|\mathcal{C}_k|$) \\
$s_i$                 & Estrategia pura actual de $i$ \\
$\mathbf{h}_i$        & Punto adelantado (look-ahead) [m] \\
$\mathrm{cnt}^{\pm}_k$& Contadores de pasos (transición de modo) \\
}

\nomblock{Parámetros del sistema}{%
$N,K$                 & Número de robots y de tareas \\
$n_k$                 & Cardinalidad mínima de la tarea $k$ \\
$T_s$                 & Período de muestreo [s] \\
$\beta$               & Pendiente de la sigmoide de payoff \\
$\lambda$             & Escala espacial del payoff [m] \\
$r_k$                 & Recompensa base de la tarea $k$ \\
$\mu_0,\mu_{\min}$    & Tasa de revisión base y piso \\
$\sigma_{\mathrm{rev}}$ & Ancho de la tasa de revisión [m] \\
$\tau_s$              & Pasos para activar el transporte \\
$\rho_s$              & Umbral de histéresis \\
$\ell$                & Distancia look-ahead [m] \\
$\Delta T_c$          & Período de replanificación T5 [s] \\
$\varepsilon$         & Ganancia de consenso \\
$\varepsilon_r$       & Regularización de repulsión [m] \\
$\varepsilon_f$       & Regularización de formación [m] \\
$r_{\mathrm{det}}$    & Radio de detección de coalición [m] \\
$r_{\mathrm{form}}$   & Radio de activación de formación [m] \\
$R$                   & Rango de comunicación [m] \\
$d^*$                 & Distancia geométrica de formación [m] \\
$\eta_o$              & Ganancia de repulsión de obstáculos \\
$\varepsilon_o$       & Regularización de obstáculos [m] \\
$c_i$                 & Capacidad del robot $i$ (ext.\ heterogénea) \\
$M_k$                 & Demanda de capacidad de la carga $k$ (ext.) \\
}

\nomblock{Funciones y operadores}{%
$\Phi(z,n)$       & Sigmoide de cardinalidad \\
$\Psi(d)$         & Descuento espacial gaussiano \\
$f_{ik}$          & Payoff de $i$ sobre la tarea $k$ \\
$F_i$             & Campo de potencial artificial \\
$V(\mathbf{x})$   & Función potencial del juego \\
$L_G,\lambda_2$   & Laplaciano de $G$ y valor de Fiedler \\
}

\nomblock{Métricas de evaluación}{%
$\Theta$                  & Throughput [min$^{-1}$] \\
$\bar{T},T_{\mathrm{coal}}$ & Tiempo de tarea y de coalición [s] \\
$F$                       & Tasa de factibilidad ($[0,1]$) \\
$D_{\mathrm{tot}}$        & Distancia total recorrida [m] \\
$D(R)$                    & Ratio de degradación $\Theta(R)/\Theta(\infty)$ \\
$\Delta J$                & Brecha de descentralización \\
$\Delta_{\mathrm{chg}}$   & Cambios de asignación por paso \\
$E_{\mathrm{DAC}}$        & Error medio de estimación de ocupación \\
$N_{\mathrm{fail}}$       & Episodios fallidos por escenario \\
}

\end{multicols}
\endgroup


\clearpage

%% ====================================================================
%% Cuerpo principal (numeración arábiga desde 1)
%% ====================================================================

\pagenumbering{arabic}
%% ====================================================================
\section{Introducción}
\label{sec:intro}
%% ====================================================================

La operación de un almacén automatizado obliga a coordinar flotas heterogéneas frente
a cargas que demandan distinto número de manipuladores. No todas pesan lo mismo. En un
escenario con quince AGV y cinco tipos de carga activos en simultáneo (la
Figura~\ref{fig:problema} ilustra un subcaso con $N{=}6$ y $K{=}3$), un solo robot basta
para cada una de tres cajas ligeras, una paleta exige dos y una estructura de acero,
cuatro. Las cargas aparecen y se completan sin pausa. La flota debe reasignarse y formar
coaliciones que satisfagan la cardinalidad mínima de cada tarea. El
resultado es una asignación distribuida con restricciones de cardinalidad que deben
cumplirse en tiempo real, sin coordinador central y con comunicación limitada al entorno
inmediato de cada robot. Difiere así de la planificación de rutas y del seguimiento de
formación fija.

Frente al transporte cooperativo clásico, donde varios robots empujan juntos una sola
carga \citep{ebel2024cooperative, rosenfelder2024force}, este problema invierte la
dificultad. Allí la composición del grupo está fijada de antemano; el reto es controlar
el conjunto. Aquí esa composición es la variable de diseño. La flota debe resolver
localmente cuántos robots asigna a cada tarea, y la decisión solo vale si produce
coaliciones factibles dentro del horizonte de operación. Los métodos de
asignación distribuida por subasta \citep{choi2009consensus} o por formación
de coaliciones \citep{dutta2021hedonic, shan2024distributed} abordan la
coordinación desde una perspectiva discreta o negociada; en la literatura revisada,
estos enfoques no se formulan como una dinámica poblacional continua acoplada a
navegación reactiva para inducir coaliciones de cardinalidad mínima. El aprendizaje
por refuerzo \citep{bezerra2025learning} produce políticas
efectivas pero dependientes de la distribución de entrenamiento y sin garantías
de robustez ante configuraciones de tarea no vistas.

\subsection{Dinámicas poblacionales como ley de control distribuida}

Una línea de investigación alternativa, desarrollada principalmente por el grupo
de Quijano, propone emplear dinámicas evolutivas de juegos poblacionales
directamente como leyes de control distribuido \citep{quijano2017population}.
La proporción de robots asignados a cada tarea se lee como el perfil de estrategias de
una población, y una dinámica de revisión actualiza ese perfil según las diferencias de
\emph{payoff} (función de utilidad que recibe cada robot por elegir una tarea). El
coordinador central desaparece. Cada agente actúa solo con información local, y el
sistema converge al equilibrio Nash bajo condiciones de juego potencial.

Como muestran \citet{barreiro2017distributed} y \citet{martinez2020formation}, las
dinámicas distributivas del replicador permiten coordinar formaciones de robots móviles
con comunicación limitada en grafo, extensión que \citet{barreiro2021uav} lleva a
formaciones de UAVs con datos reales. Más recientemente, \citet{dinh2026bsplines} proponen
mecanismos basados en B-splines para diseñar funciones de payoff que satisfagan
restricciones de formación específicas. En esa línea, sin embargo, suele suponerse que
el tamaño del grupo que realiza cada tarea es fijo y conocido de antemano. El requisito
de cardinalidad mínima dinámica no recibe tratamiento explícito.

\subsection{El papel de la política de revisión: replicador y Smith}
\label{subsec:porquesmith}

Manteniendo fija la arquitectura de estimación, una de las preguntas que la escalera
comparativa del trabajo formaliza es qué aporta la \emph{política de revisión} de una
dinámica poblacional. Los tratamientos T3 (replicador distribuido)
y T4 (dinámicas de Smith) comparten el mismo payoff y el mismo estimador por
consenso, y difieren únicamente en esta ley; comparar ambos aísla su efecto. Ambas
parten del mismo marco. El estado de la población es el vector de fracciones
$\mathbf{x} = (x_k)$, donde $x_k$ es la proporción de robots asignados a la tarea $k$.
Cada robot estima $x_k$ por consenso promedio distribuido y obtiene así
$\hat{x}_{ik}$. Hasta aquí, Smith y DRD son idénticos: ambos requieren esta
\emph{primera capa de consenso}.

La diferencia aparece en lo que cada dinámica requiere más allá de esa primera capa.
Las dinámicas del replicador (DRD) comparan el payoff de cada estrategia con la
\emph{media poblacional} $\bar{f}(\mathbf{x}) = \sum_k x_k f_k(\mathbf{x})$:
esta cantidad es global porque depende de la distribución completa $\mathbf{x}$,
no solo del estado local de un agente. Estimarla distribuida\-mente exige una
\emph{segunda capa de consenso} sobre los payoffs individuales $f_k(\hat{x}_i)$.
Cuando esta segunda capa opera sobre un grafo irregular que cambia con el movimiento
de los robots, el operador de consenso actúa como un filtro pasa-bajos espectral
\citep{sandryhaila2013dsp}: atenúa las componentes de alta frecuencia del vector
de payoffs sobre la estructura del Laplaciano de $G(t)$ e introduce un sesgo
dependiente de la topología $\delta(G, \beta)$ que desplaza el equilibrio efectivo
respecto al equilibrio Nash verdadero.

A diferencia del replicador, las dinámicas de Smith evitan esta segunda capa porque la
actualización se basa en comparaciones bilaterales $\pos{f_{ik} - f_{ij}}$ calculadas localmente con los
payoffs propios: ningún agente necesita estimar la media poblacional global. La
condición de punto fijo $\pos{f_{ik} - f_{ij}} = 0$ se satisface exactamente con
información local derivada de $\hat{x}_{ik}$. Además, las estrategias extintas
($p_{ik} = 0$) pueden recuperarse, ya que la tasa de migración es independiente
de la masa actual en cada estrategia. La línea base T3 (tratamiento de replicador
distribuido) cuantifica empíricamente el efecto de este sesgo bajo las mismas
topologías de comunicación que el método propuesto.

\subsection{La contribución de este trabajo}

La pieza central del trabajo es una función de payoff con componente sigmoide decreciente
que representa requisitos mínimos de cardinalidad mediante cotas inferiores, en contraste
con las cotas superiores de \citet{martinezpiazuelo2022automatica}. En el caso base con
$\Psi\equiv1$ preserva una estructura de juego de congestión ponderado. Al activar
$\Psi(d)$, el sistema implementado se vuelve una perturbación espacial que acopla
asignación y movimiento. Sobre ese payoff se adaptan y comparan el replicador distribuido
y las dinámicas de Smith bajo una misma arquitectura de estimación por consenso dinámico
de media; así se aísla el efecto de la política de revisión. Una tasa de revisión
espacialmente modulada suprime, además, las reasignaciones espurias durante la
aproximación a las cargas y en régimen estacionario. Queda un último paso: medir el coste
de la descentralización frente a una referencia centralizada replanificada y la
degradación al reducir el rango de comunicación. Estos cuatro elementos definen una
arquitectura donde la estimación de ocupación, la revisión estratégica y la navegación
reactiva se actualizan en un mismo bucle temporal, sin coordinador central de asignación.

La validación se realiza mediante simulación cuantitativa en Python con seis
escenarios y cinco tratamientos comparativos dispuestos como escalera incremental
(T1 heurística voraz local, T2 DAC con regla voraz, T3 replicador distribuido
modificado, T4 dinámicas de Smith con revisión modulada (método propuesto) y
T5 referencia centralizada replanificada), más una validación cualitativa y
cinemática en CoppeliaSim.

\subsection{Estructura del documento}

La sección~\ref{sec:planteamiento} describe el problema de asignación distribuida
con restricciones de cardinalidad y la brecha que lo motiva. La
sección~\ref{sec:hipotesis} formula la pregunta de investigación y las hipótesis.
Los objetivos se especifican en la sección~\ref{sec:objetivos}, y el alcance y
limitaciones en la sección~\ref{sec:alcance}. La sección~\ref{sec:metodologia},
núcleo de esta entrega, desarrolla las tres dinámicas acopladas, la
función de payoff, el análisis de convergencia y el diseño experimental completo.
Las secciones \ref{sec:marco}--\ref{sec:conclusiones} están reservadas para el
marco teórico, resultados y conclusiones de la entrega final.

\subsection{Estado actual del TFM}
\label{subsec:estado}

A la fecha de esta entrega (mayo de 2026) se han completado las siguientes etapas:

\begin{itemize}
  \item \textbf{Revisión bibliográfica} (21 referencias en 6 áreas temáticas:
    dinámicas poblacionales, consenso distribuido, formación de coaliciones,
    transporte cooperativo, juegos potenciales y procesado de señales en grafos).
  \item \textbf{Formulación matemática completa}: modelo de robots, tareas y
    grafo de comunicación; función de payoff multiplicativa con sigmoide
    decreciente; tres dinámicas acopladas (consenso, Smith en tiempo discreto,
    gradiente reactivo); análisis formal de convergencia para el caso simétrico
    (Teorema~\ref{prop:convergencia} con demostración y cálculo de Hessiana).
  \item \textbf{Diseño experimental}: seis escenarios, cinco tratamientos
    comparativos (incluido el replicador distribuido como línea base clave),
    once métricas definidas y criterios de aceptación conjuntivos.
  \item \textbf{Pseudocódigo formal} del bucle de control distribuido
    (Algoritmo~\ref{alg:loop}).
\end{itemize}

En curso: implementación en Python de las tres dinámicas acopladas, los
tratamientos comparativos y la campaña experimental. El plan de trabajo
completo se detalla en la Sección~\ref{sec:cronograma}.

\subsection{Planteamiento del problema y justificación}
\label{sec:planteamiento}

\subsubsection{Descripción formal del problema}

El Anexo~I de este trabajo mencionaba ``cargas heterogéneas'' (objetos que difieren
en peso, fragilidad y tamaño). Aquí esa heterogeneidad se formaliza
como un requisito mínimo de cardinalidad $n_k$: el número de manipuladores
que una carga determinada exige para ser transportada con seguridad. La abstracción es
intencional: atributos como la masa o la fragilidad se condensan en un valor escalar
$n_k$ que el sistema de asignación consume; derivar $n_k$ a partir de un modelo de contacto explícito (cierre de
fuerza, estabilidad del agarre) es una capa complementaria documentada en
\citet{ebel2024cooperative} y \citet{rosenfelder2024force} que queda fuera del
alcance de este TFM (véase \S\,\ref{sec:alcance}). El problema central y el objeto de
estudio del Anexo~I se mantienen; este informe concreta la terminología (formación de
coaliciones) y la metodología adoptada.

Con esta abstracción, el sistema se define así. Se considera
una flota de $N$ AGV que operan en un entorno 2D. El conjunto de robots
se indexa $i \in \{1, \ldots, N\}$ y el conjunto de tareas activas en un instante
dado se indexa $k \in \{1, \ldots, K\}$. Cada tarea $k$ está ubicada en una
posición conocida $\mathbf{l}_k \in \RR^2$ y exige que al menos $n_k \in \NN$
robots lleguen a ella para que la carga pueda ser manipulada. La tarea $k=0$ representa
el estado de inactividad: un robot inactivo (\textit{idle}) no contribuye a ninguna tarea.

Cada robot mantiene un vector de estrategias $\mathbf{p}_i(t) \in \simplex^K$,
donde $\simplex^K = \{p \in \RR^{K+1} : \sum_{k=0}^K p_k = 1,\, p_k \geq 0\}$
es el símplex de probabilidad. El componente $p_{ik}(t)$ representa la propensión
del robot $i$ a seleccionar la tarea $k$. La estrategia efectiva del robot es
$s_i(t) = \argmax_k p_{ik}(t)$. El objetivo del sistema es que, para cada tarea $k$,
la coalición $\mathcal{C}_k(t) = \{i : s_i(t) = k, \norm{\mathbf{q}_i - \mathbf{l}_k}
< r_{\text{det}}\}$ satisfaga $|\mathcal{C}_k(t)| \geq n_k$ en tiempo finito, y
que esta condición se mantenga robustamente frente a cambios en el conjunto de
tareas activas.

\begin{figure}[htbp]
\centering
\begin{tikzpicture}[x=1cm, y=1cm,
  agv/.style={circle, draw=VIUOrange, fill=VIUOrange!20,
    minimum size=0.62cm, font=\scriptsize\bfseries, inner sep=0pt},
  pickup/.style={rectangle, draw=VIUDark, fill=VIUDark!8,
    minimum width=0.88cm, minimum height=0.88cm, font=\tiny, align=center, inner sep=2pt},
  comm/.style={draw=VIUGray!70, thin, dashed},
  motion/.style={-{Stealth[length=4pt]}, draw=VIUOrange!75, thick, shorten >=3pt, shorten <=3pt}
]
%-- Planta de bodega
\draw[VIUDark!40, thick, rounded corners=5pt] (0,0) rectangle (11.5,7.8);
\node[font=\tiny, VIUGray, anchor=north west] at (0.15,7.65)
  {Almacén automatizado ($20\,\mathrm{m}\times 14\,\mathrm{m}$)};

%-- Tareas: posiciones de recogida l_k
\node[pickup] (L1) at (2.0,6.5) {$k{=}1$\\$n_1{=}1$};
\node[pickup] (L2) at (5.5,3.2) {$k{=}2$\\$n_2{=}3$};
\node[pickup] (L3) at (9.5,6.0) {$k{=}3$\\$n_3{=}2$};
\node[font=\tiny, VIUDark, above=1pt of L1] {$\mathbf{l}_1$};
\node[font=\tiny, VIUDark, above=1pt of L2] {$\mathbf{l}_2$};
\node[font=\tiny, VIUDark, above=1pt of L3] {$\mathbf{l}_3$};

%-- Destinos d_k (dibujados como rombos manuales)
\def\rdia{0.30}
\foreach \cx/\cy/\lbl in {1.5/1.0/$\mathbf{d}_1$, 8.5/0.8/$\mathbf{d}_2$, 10.8/3.8/$\mathbf{d}_3$}{%
  \draw[VIUDeep, thick]
    (\cx,\cy+\rdia) -- (\cx+\rdia,\cy) -- (\cx,\cy-\rdia) -- (\cx-\rdia,\cy) -- cycle;
  \node[font=\tiny, VIUDeep] at (\cx,\cy) {\lbl};
}

%-- Robots
\node[agv] (R1) at (2.0,5.0)  {1};
\node[agv] (R2) at (4.4,2.7)  {2};
\node[agv] (R3) at (6.5,2.7)  {3};
\node[agv] (R4) at (5.6,4.2)  {4};
\node[agv] (R5) at (8.5,5.0)  {5};
\node[agv] (R6) at (9.8,4.8)  {6};

%-- Coalicion C_2 = {2,3,4}: recuadro resaltado
\draw[VIUOrange, thick, dashed, rounded corners=6pt, fill=VIUOrange!5]
  (3.7,2.0) rectangle (7.2,5.0);
\node[font=\tiny, VIUOrange, anchor=south west] at (3.7,5.0)
  {$\mathcal{C}_2$: $|\mathcal{C}_2|=3\geq n_2$ $\checkmark$};

%-- Flechas de transporte: coalicion 2 se mueve hacia d_2
\draw[{Stealth[length=4pt]}-{Stealth[length=4pt]}, VIUOrange!60, thick, dashed]
  ($(L2.south)+(0,-0.1)$) -- (8.5,0.8)
  node[midway, below, font=\tiny, VIUGray] {transporte $\rightarrow \mathbf{d}_2$};

%-- Links de comunicacion dentro de la coalicion
\draw[comm] (R2) -- (R3);
\draw[comm] (R2) -- (R4);
\draw[comm] (R3) -- (R4);

%-- Robots 5 y 6: convergiendo hacia L3
\draw[comm] (R5) -- (R6);
\draw[motion] (R5) -- ($(L3.south west)+(0.05,0.05)$);
\draw[motion] (R6) -- ($(L3.east)+(-0.05,0)$);

%-- Robot 1: solo, hacia L1
\draw[motion] (R1) -- ($(L1.south)+(0,0.05)$);

%-- Circulo de rango R (ejemplo alrededor de R5)
\draw[VIUGray, dotted, thin] (R5) circle (1.4cm);
\node[font=\tiny, VIUGray, anchor=north] at ($(R5)+(0,-1.4)$) {rango $R$};

%-- Leyenda
\begin{scope}[shift={(0.1,0.05)}]
  \node[agv, minimum size=0.42cm] at (0.38,0.38) {};
  \node[font=\tiny, anchor=west] at (0.65,0.38) {AGV $i$};
  \node[pickup, minimum width=0.42cm, minimum height=0.42cm] at (2.5,0.38) {};
  \node[font=\tiny, anchor=west] at (2.8,0.38) {Recogida $\mathbf{l}_k$};
  \draw[VIUDeep, thick]
    (5.2,0.62) -- (5.45,0.38) -- (5.2,0.14) -- (4.95,0.38) -- cycle;
  \node[font=\tiny, anchor=west] at (5.55,0.38) {Entrega $\mathbf{d}_k$};
  \draw[comm] (7.3,0.38) -- (8.0,0.38);
  \node[font=\tiny, anchor=west] at (8.05,0.38) {Enlace com.};
\end{scope}
\end{tikzpicture}
\caption[Escenario de asignación distribuida con coaliciones]%
{Escenario del problema de asignación distribuida con coaliciones. Seis AGV
($N=6$) en un almacén automatizado con tres tareas activas ($K=3$). La coalición $\mathcal{C}_2$
(robots~2, 3, 4) ya satisface $|\mathcal{C}_2|=3=n_2$ y transita hacia el destino
$\mathbf{d}_2$. Los robots~5 y~6 convergen hacia la tarea~3 ($n_3=2$). El robot~1 atiende
solo la tarea~1 ($n_1=1$). El círculo punteado indica el rango de comunicación~$R$.}
\label{fig:problema}
\end{figure}


Dos dificultades distinguen este problema de la asignación estática de tareas. La
restricción $|\mathcal{C}_k| \geq n_k$ es una \emph{cota inferior} dinámica: si tres
robots abandonan una tarea que requiere cuatro, el sistema debe reclutar uno nuevo sin
coordinación central. A ello se suma que, cuando las tareas aparecen, desaparecen o
cambian su prioridad durante la operación, las coaliciones deben reconstituirse sin
reinicializar el sistema completo.

\subsubsection{Límites de las soluciones existentes}

Los métodos de asignación descentralizada basados en subastas
\citep{choi2009consensus} y en formación de coaliciones hedo\-nistas
\citep{dutta2021hedonic, zhang2024coalition} coordinan la asignación mediante
negociación iterativa entre agentes, pero producen asignaciones discretas que
no se actualizan continuamente ante cambios en el entorno. El aprendizaje por refuerzo
\citep{bezerra2025learning} genera políticas efectivas en distribución, pero exige
entrenamiento previo sobre escenarios representativos y carece de garantías formales de
convergencia ante configuraciones de tarea no vistas. Ambas familias comparten un mismo
límite. No proporcionan una ley de control diferencial que actualice sin interrupción
las estrategias de asignación en tiempo real.

En cambio, las dinámicas poblacionales \citep{barreiro2021uav, dinh2026bsplines,
hurtado2024potential} sí ofrecen leyes de control con garantías de convergencia
al equilibrio Nash, gracias a la estructura de juego potencial subyacente. Aun así,
en todos los trabajos de esta línea el tamaño del grupo que realiza
cada tarea es un parámetro del diseño fijado de antemano: los robots se asignan
en fracciones fijas, no en coaliciones cuyo tamaño emerge de la dinámica.
Cuando el tamaño de la coalición pasa a ser la variable de estado (porque cada carga
impone su propio requisito mínimo), las dinámicas deben reclutar o liberar
robots según las necesidades de cada carga; ninguna formulación con ese grado de
generalidad apareció en los trabajos analizados.

El trabajo de \citet{martinezpiazuelo2022automatica} introduce restricciones de
capacidad en juegos poblacionales, pero estas restricciones son \emph{cotas superiores}
sobre cuántos agentes pueden elegir una estrategia: impiden que una estrategia se
sature por encima de un umbral. El caso que aquí se aborda requiere cotas inferiores:
la tarea no puede ejecutarse si hay \emph{menos} de $n_k$ robots. El cambio de signo no
es cosmético. Los resultados de estabilidad de ese trabajo no se trasladan directamente
a este escenario.

En las referencias consultadas tampoco se identifica un enfoque que trate
explícitamente el caso $N < \sum_k n_k$ (menos robots que los requeridos en total):
en ese régimen la
flota no puede satisfacer todos los requisitos de cardinalidad simultáneamente y
el sistema debe priorizar. Este caso límite, relevante en escenarios con alta
demanda transitoria, queda fuera del alcance de este trabajo pero se identifica como
línea futura.

\subsubsection{La brecha específica}

La combinación de tres elementos (dinámicas de Smith como ley de control directa,
función de payoff con cota inferior de cardinalidad mediante sigmoide decreciente
y tasa de revisión espacialmente modulada) no aparece en el corpus
bibliográfico revisado. Cada elemento por separado tiene antecedentes:
las dinámicas de Smith en sistemas de control \citep{barreiro2017distributed},
las funciones sigmoide para incentivos de reclutamiento en antropología
e inteligencia de enjambre \citep{theraulaz1999stigmergy},
y las tasas de revisión moduladas por contexto en juegos evolutivos
\citep{sandholm2010population}. La integración de los tres, con el análisis
formal que ella requiere, constituye la contribución de este trabajo.

La afirmación del trabajo es acotada a propósito: bajo comunicación local y
cargas heterogéneas con requisitos mínimos de cardinalidad, ¿puede una arquitectura
distribuida e interpretable que integre en un único bucle la estimación de ocupación,
la revisión poblacional de estrategias y la navegación reactiva formar coaliciones
factibles con una pérdida de desempeño acotada frente a una referencia centralizada?
Responderla es el objetivo. Afirmar la superioridad del control distribuido sobre el
centralizado en todo régimen operativo queda fuera de su alcance.



%% ====================================================================
\section{Objetivos}
\label{sec:objetivos}
%% ====================================================================

A partir del problema planteado en la introducción, esta sección formula el objetivo
general del TFM y lo desglosa en objetivos específicos, medibles y alcanzables. Los
objetivos específicos guían el desarrollo metodológico y la campaña experimental, y su
grado de cumplimiento se contrastará en la entrega final.

\subsection{Objetivo general}

Diseñar y evaluar, mediante simulación cuantitativa reproducible, un esquema de
coordinación distribuida local basado en consenso dinámico de media y dinámicas
poblacionales para la formación adaptativa de coaliciones en flotas multi-AGV que
transportan cooperativamente cargas heterogéneas con requisitos mínimos de
cardinalidad, integrando estimación de ocupación, asignación a cargas, formación de
coaliciones y navegación reactiva en un único bucle de control, caracterizado mediante
una escalera comparativa de cinco tratamientos (T1--T5) y contrastado con una
referencia centralizada replanificada bajo seis escenarios experimentales.

\subsection{Objetivos específicos}

\begin{enumerate}
  \item Formalizar el problema de formación de coaliciones multi-AGV como un
    sistema de asignación distribuida con requisitos mínimos de cardinalidad,
    estrategia de inactividad, comunicación local y navegación cinemática en un
    entorno 2D.

  \item Diseñar una función de payoff con componente sigmoide decreciente de
    cardinalidad y descuento espacial, distinguiendo entre el juego base con
    $\Psi\equiv1$ (donde se analiza la estructura de juego potencial y se caracteriza
    el intervalo $(\beta_{\min}, \beta_{\max})$ de estabilidad local del equilibrio
    mediante un análisis de Lyapunov) y el sistema implementado con $\Psi(d)$,
    tratado como una perturbación espacial validada por simulación.

  \item Proponer la tasa de revisión espacialmente modulada $\mu_i(t)$ y
    cuantificar su efecto sobre la frecuencia de reasignaciones y el tiempo de
    formación de coaliciones en los experimentos.

  \item Implementar las tres dinámicas acopladas (consenso distribuido,
    Smith en tiempo discreto y gradiente reactivo con evitación de obstáculos
    estáticos) en un paquete Python
    reproducible, con soporte para cinco tratamientos y seis escenarios con semillas
    controladas y registro completo de trayectorias.

  \item Validar cualitativamente el comportamiento de formación de coaliciones y el
    transporte cooperativo de la coalición con evitación reactiva de obstáculos en
    CoppeliaSim con robots Pioneer~3-DX.
\end{enumerate}


%% ====================================================================
\section{Hipótesis de partida}
\label{sec:hipotesis}
%% ====================================================================

\textbf{Pregunta de investigación.}
?`Pueden las dinámicas de Smith con una función de payoff multiplicativa y cota
inferior sigmoide de cardinalidad garantizar, de forma distribuida y sin coordinador
central, que una flota de AGV forme coaliciones de tamaño mínimo requerido para
el transporte cooperativo de cargas, y cómo cambia este comportamiento al variar
el rango de comunicación?

\textbf{Hipótesis~1.}
Bajo las dinámicas de Smith con la función de payoff propuesta, el sistema
converge a un equilibrio Nash en el que cada tarea $k$ tiene exactamente $n_k$ robots
asignados, en el caso simétrico balanceado ($N = Kn$, $n_k = n$, $\Psi \equiv 1$,
$R = \infty$). La condición suficiente de estabilidad local exige que la pendiente
$\beta$ de la función sigmoide pertenezca al intervalo $(\beta_{\min}, \beta_{\max})$.

\textbf{Hipótesis~2.}
La tasa de revisión espacialmente modulada $\mu_i(t)$ reduce significativamente
el número de cambios de asignación por robot durante el estado estacionario
respecto a una tasa constante $\mu_i = \mu_0$, sin degradar de forma relevante
el tiempo de convergencia de la coalición.

\textbf{Hipótesis~3.}
El método propuesto supera en tasa de factibilidad de formación de coaliciones
a la heurística voraz local y al consenso puro en todos los escenarios evaluados,
y la ventaja se mantiene cuando el rango de comunicación $R$ se reduce,
aunque la brecha disminuye con la conectividad.

\textbf{Hipótesis~4.}
La brecha de descentralización $\Delta J = (\Theta_{\text{cent}} -
\Theta_{\text{dist}}) / \max(\Theta_{\text{cent}}, \delta_0)$, con $\delta_0 = 10^{-3}$,
permanece acotada en los escenarios homogéneos y heterogéneos, con valores
aceptables para un entorno de almacén con comunicación imperfecta.

\subsection{Criterios cuantitativos de aceptación}
\label{subsec:criterios}

Los umbrales de aceptación que operacionalizan las hipótesis son los siguientes:

\begin{description}
  \item[H1 (convergencia Nash):] El número de episodios con convergencia
    ($|x_k(T_{\text{conv}}) - x_k^*| < 0{,}01$ para todo $k$) es $\geq 27$ de 30
    repeticiones en el Escenario~A balanceado ($N = Kn = 6$, $n_k = 2$).

  \item[H2 (revisión modulada):] $\Delta_{\text{chg}}$ del tratamiento T5 con
    tasa espacial es al menos un 40\,\% menor que con tasa constante $\mu_i = \mu_0$,
    sin que $T_{\text{coal}}$ aumente más de un 15\,\%. El 40\,\% se justifica
    porque, al valor nominal $\sigma_{\text{rev}} = 2\,\text{m}$, un robot a
    distancia $d = \sigma_{\text{rev}}$ tiene $\mu_i \approx 0{,}39\,\mu_0$; en
    estado estacionario la mayoría de robots están dentro de esa distancia,
    por lo que una reducción agregada del 40\,\% es conservadora.

  \item[H3 (factibilidad de coalición):] La tasa de factibilidad $F \geq 0{,}85$
    para T5 en todos los escenarios evaluados, incluyendo $R = 3\,\text{m}$.
    El umbral del 85\,\% refleja que una tasa de fallo del 15\,\% es recuperable
    mediante reintento en entornos reales; umbrales más laxos no serían informativos.

  \item[H4 (brecha de descentralización):] $\Delta J < 0{,}25$ en escenarios
    homogéneos y $\Delta J < 0{,}40$ en heterogéneos. Una brecha del 25\,\%
    significa que el sistema distribuido completa al menos 3 tareas de cada 4 que
    completaría un planificador centralizado con estado global perfecto.
\end{description}

Las hipótesis se contrastan mediante campaña de simulación con $m \geq 30$
repeticiones por tratamiento-escenario, bajo las mismas semillas, trayectorias
y parámetros iniciales para todos los métodos. La Hipótesis~1 se verifica
numéricamente en el escenario~A antes de ejecutar la campaña completa.


%% ====================================================================
\section{Metodología}
\label{sec:metodologia}
%% ====================================================================

En este capítulo se delimita primero el alcance y las limitaciones del estudio y se
sitúa el plan de trabajo; a continuación se formaliza el modelo del sistema y se
describen las tres dinámicas acopladas que conforman el método. El capítulo cierra
con el análisis de convergencia, el diseño experimental, las métricas y los criterios
de aceptación.

\subsection{Alcance y limitaciones}
\label{sec:alcance}

Esta sección delimita qué aspectos del problema aborda el TFM y cuáles quedan fuera, y
sitúa el trabajo realizado dentro del cronograma global. La distinción explícita entre
alcance y limitaciones acota las afirmaciones que el trabajo puede sostener y orienta
la interpretación de los resultados de la entrega final.

\subsubsection{Alcance}

El trabajo se limita a un entorno 2D con robots de cinemática uniciclo. La validación
es exclusivamente por simulación: Python proporciona los resultados cuantitativos y
CoppeliaSim proporciona la validación visual aplicada. El sistema puede manejar hasta
$N = 15$ AGV y $K = 5$ tareas activas simultáneas; el límite $N = 15$ responde
a la complejidad computacional del simulador Python en tiempo real y al rango de
flotas AGV típico en almacenes medianos, no a una restricción teórica del método. El rango de comunicación $R$
se trata como parámetro variable entre 1\,m e infinito. Las cargas se modelan como puntos de destino con requisito mínimo de cardinalidad
$n_k$. Esta es una decisión de \textbf{descomposición por capas}: la capa de
asignación (este TFM) determina \emph{qué} robots forman la coalición y
\emph{cuándo} inician el transporte; la capa de manipulación (no modelada aquí)
determina \emph{cómo} aplican fuerzas de contacto para mover la carga de forma
estable. Esta segunda capa está documentada en \citet{ebel2024cooperative} y
\citet{rosenfelder2024force}, que resuelven el control de fuerzas para empuje y
agarre no-prehensil con robots diferenciales. No se modela el contacto mecánico
entre robot y carga durante el transporte; la fase de transporte propaga
cinemáticamente el centroide de la coalición hacia $\mathbf{d}_k$.

El trabajo incluye explícitamente las tres dinámicas acopladas ejecutadas en un
único bucle de control cada $T_s = 0{,}1\,\text{s}$, la función de payoff
multiplicativa con parámetros $\beta$, $\lambda$, $r_k$, y el análisis formal
de convergencia para el caso simétrico con comunicación global
(Teorema~\ref{prop:convergencia}). El diseño experimental contempla seis
escenarios con cinco tratamientos comparativos y al menos 30 repeticiones por
combinación, validación cualitativa en CoppeliaSim con robots Pioneer~3-DX
(Escenario~F), y un paquete Python documentado y reproducible con semillas
fijadas, configuraciones YAML y registro completo de trayectorias.

El transporte de la coalición incorpora evitación reactiva de obstáculos estáticos a
nivel cinemático del centroide de la coalición, manteniendo la geometría de formación
$d^*$. Esta capa no modela la dinámica de contacto entre robots y carga ni ofrece
garantías formales de evitación de colisiones: la navegación con obstáculos se evalúa
exclusivamente de forma empírica.

\subsubsection{Limitaciones}

Quedan fuera del alcance la experimentación con hardware real, la percepción
visual, SLAM, evitación de obstáculos complejos y planificación global de rutas.
El análisis formal de convergencia (Teorema~\ref{prop:convergencia}) se
restringe al caso simétrico con comunicación global ($R = \infty$) y tareas
homogéneas; la extensión al caso general con grafo irregular y tareas
heterogéneas queda como trabajo futuro. Esa extensión conecta con el
formalismo de dinámicas del replicador distribuido con procesado en grafos
(DRD-PI), donde la convolución de payoffs sobre la estructura espectral del
Laplaciano de $G(t)$ sigue la teoría de procesado de señales en grafos de
\citet{sandryhaila2013dsp}. Tampoco se aborda la generación dinámica de tareas
durante la ejecución, aunque los experimentos de perturbación (Escenario~E)
prueban la reactividad ante cambios en el conjunto de tareas.

La cota inferior de cardinalidad introduce rigidez numérica cuando $\beta$ es
muy grande (la función sigmoide se aproxima a una función escalón). Esta zona
de parámetros se excluye explícitamente del análisis: el TFM opera en
$\beta \in (\beta_{\min}, \beta_{\max})$. Incluso dentro de ese intervalo, en
escenarios donde el número total de robots es inferior a la demanda total
($N < \sum_k n_k$), el sistema no puede satisfacer todos los requisitos de
cardinalidad simultáneamente: las dinámicas de Smith oscilarán sin alcanzar
el equilibrio interior y el estado estacionario dependerá de las prioridades
relativas de las tareas. Este régimen queda fuera del alcance del análisis
formal. Además, la tasa de revisión modulada reduce pero no elimina las
oscilaciones de alta frecuencia en la frontera del símplex (\emph{chattering})
cuando $\beta$ está cercano a $\beta_{\max}$; este comportamiento es observable
en simulación y se presentará en los resultados de la entrega final.

La evitación de obstáculos se implementa mediante un término repulsivo de campo de
potencial; al heredar las limitaciones clásicas de los campos de potencial artificial
(mínimos locales), no se garantiza la ausencia de atascos ni la completitud de la
navegación. Los obstáculos se modelan como elementos estáticos convexos simples; no se
aborda planificación global de rutas ni SLAM.

El núcleo formal y experimental supone flota homogénea (cardinalidad mínima). La
heterogeneidad de capacidad entre robots se formula como extensión
(\S\,\ref{subsec:heterogenea}) y se explora cualitativamente (Escenario~G), sin
garantía formal.

\subsubsection{Plan de trabajo y cronograma}
\label{sec:cronograma}

El cronograma (Figura~\ref{fig:gantt}) abarca desde el inicio oficial del TFM
(4 de marzo de 2026) hasta la primera convocatoria de defensa (21--25 de
septiembre de 2026), con tres tutorías colectivas síncronas en los hitos T0,
T1 y T2. Las etapas completadas aparecen en naranja; las pendientes, en gris claro.

\begin{figure}[htbp]
\centering
\begin{tikzpicture}[x=1.55cm, y=0.72cm]

\tikzset{
  done/.style    = {fill=VIUOrange!70, draw=VIUDeep,    line width=0.45pt, rounded corners=2pt},
  current/.style = {fill=VIUOrange,    draw=VIUDeep,    line width=0.75pt, rounded corners=2pt},
  pending/.style = {fill=VIUGray!12,   draw=VIUGray!45, line width=0.45pt, rounded corners=2pt},
  defense/.style = {fill=VIUDeep!18,   draw=VIUDeep,    line width=0.75pt, rounded corners=2pt}
}
\def\BH{0.50}

%-- Franjas alternadas de fondo (cada dos meses para guiar la lectura)
\foreach \m in {0,2,4,6}{
  \fill[VIUGray!7] (\m, 0.68) rectangle (\m+1, 7.58);
}
%-- Separadores de mes
\foreach \m in {0,...,7}{
  \draw[VIUGray!30, very thin] (\m, 0.68) -- (\m, 7.58);
}
%-- Línea de base horizontal
\draw[VIUGray!35, thin] (0, 0.68) -- (7, 0.68);

%-- Etiquetas de meses (encabezado superior)
\foreach \m/\lbl in {0/Mar,1/Abr,2/May,3/Jun,4/Jul,5/Ago,6/Sep}{
  \node[font=\footnotesize\bfseries, color=VIUGray!80!black, anchor=south]
    at (\m+0.5, 7.60) {\lbl};
}
\node[font=\small\bfseries, color=VIUDark, anchor=south] at (3.5, 8.10) {2026};

%-------------------------------------------------------------
%  BARRAS DE FASES + etiquetas (derecha de la barra, izquierda del diagrama)
%-------------------------------------------------------------

% E1 — Revisión bibliográfica  [Mar 4 → May]  COMPLETADO
\fill[done] (0.10, 7.0) rectangle (2.0, 7.0+\BH);
\node[font=\footnotesize, color=VIUDark, anchor=east, text width=3.7cm, align=right]
  at (-0.10, 7.25) {Revisión bibliográfica};

% E2 — Formulación del problema  [Mar 4 → 29 may]  ENTREGADO (esta Tarea)
\fill[current] (0.10, 6.0) rectangle (2.9, 6.0+\BH);
\node[font=\footnotesize, color=VIUDark, anchor=east, text width=3.7cm, align=right]
  at (-0.10, 6.25) {Formulación del problema};

% E3 — Implementación Python  [Jun → mid-ago]  PENDIENTE
\fill[pending] (2.0, 5.0) rectangle (4.5, 5.0+\BH);
\node[font=\footnotesize, color=VIUDark, anchor=east, text width=3.7cm, align=right]
  at (-0.10, 5.25) {Implementación Python};

% E4 — Simulación experimental  [Jun → ago]  PENDIENTE
\fill[pending] (3.1, 4.0) rectangle (5.5, 4.0+\BH);
\node[font=\footnotesize, color=VIUDark, anchor=east, text width=3.7cm, align=right]
  at (-0.10, 4.25) {Simulación experimental};

% E5 — Validación en CoppeliaSim  [mid-ago → sep]  PENDIENTE
\fill[pending] (4.5, 3.0) rectangle (6.0, 3.0+\BH);
\node[font=\footnotesize, color=VIUDark, anchor=east, text width=3.7cm, align=right]
  at (-0.10, 3.25) {Validación CoppeliaSim};

% E6 — Redacción de la memoria  [jul → sep]  PENDIENTE
\fill[pending] (4.0, 2.0) rectangle (6.7, 2.0+\BH);
\node[font=\footnotesize, color=VIUDark, anchor=east, text width=3.7cm, align=right]
  at (-0.10, 2.25) {Redacción de la memoria};

% E7 — Defensa TFM  [21–25 sep]  PENDIENTE
\fill[defense] (6.5, 1.0) rectangle (7.0, 1.0+\BH);
\node[font=\footnotesize, color=VIUDark, anchor=east, text width=3.7cm, align=right]
  at (-0.10, 1.25) {Defensa TFM};

%-------------------------------------------------------------
%  TUTORÍAS COLECTIVAS  (triángulos sobre el límite superior)
%  T0 = 4 mar (x≈0.10)   T1 = 3 jun (x≈3.08)   T2 = 15 jul (x≈4.47)
%-------------------------------------------------------------
\foreach \tx in {0.10, 3.08, 4.47}{
  \fill[VIUGray!55]
    (\tx, 7.58) -- (\tx-0.08, 7.74) -- (\tx+0.08, 7.74) -- cycle;
}
\node[font=\tiny\bfseries, color=VIUGray!80, anchor=south west] at (0.18,  7.74) {T0};
\node[font=\tiny\bfseries, color=VIUGray!80, anchor=south west] at (3.16,  7.74) {T1};
\node[font=\tiny\bfseries, color=VIUGray!80, anchor=south west] at (4.55,  7.74) {T2};

%-------------------------------------------------------------
%  HITOS
%-------------------------------------------------------------

% Tarea 1 / Informe Intermedio  (29 may → x = 2.9)
\draw[VIUDeep, line width=1.5pt] (2.9, 0.68) -- (2.9, 7.58);
\fill[VIUDeep] (2.9, 0.68) circle(3.2pt);
\node[font=\tiny\bfseries, color=VIUDeep, anchor=north, align=center]
  at (2.9, 0.62) {Tarea\,1\\29~may.};

% Defensa, 1\textordfeminine{} convocatoria  (21 sep → x = 6.67)
\draw[VIUDark!55, line width=1.0pt, dashed] (6.67, 0.68) -- (6.67, 7.58);
\fill[VIUDark!55] (6.67, 0.68) circle(2.5pt);
\node[font=\tiny, color=VIUDark!70, anchor=north east, align=right]
  at (6.67, 0.62) {Defensa\\21--25~sep.};

%-------------------------------------------------------------
%  LEYENDA
%-------------------------------------------------------------
\fill[done]    (0.00,-0.72) rectangle (0.90,-0.35);
\node[font=\footnotesize, color=VIUDark, anchor=west] at (1.00,-0.535) {Completado};
\fill[current] (2.60,-0.72) rectangle (3.50,-0.35);
\node[font=\footnotesize, color=VIUDark, anchor=west] at (3.60,-0.535) {Entregado};
\fill[pending] (5.10,-0.72) rectangle (6.00,-0.35);
\node[font=\footnotesize, color=VIUDark, anchor=west] at (6.10,-0.535) {Pendiente};

\end{tikzpicture}
\caption[Cronograma del TFM]%
{Cronograma del TFM (marzo--septiembre 2026).
Tutorías colectivas síncronas: T0~inicio (4~mar.), T1~intermedia (3~jun.),
T2~final (15~jul.).
Fases: E1~Revisión bibliográfica, E2~Formulación del problema (Tarea~1),
E3~Implementación Python, E4~Simulación experimental,
E5~Validación CoppeliaSim, E6~Redacción de la memoria,
E7~Defensa (1.\textordfeminine\ conv., 21--25~sep.\ 2026).
La línea sólida marca el hito actual (29~may.\ 2026); la de trazos,
la convocatoria de defensa.}
\label{fig:gantt}
\end{figure}



%% ====================================================================

La investigación adopta un diseño cuantitativo, experimental en simulación y
comparativo. Más que demostrar la superioridad de una técnica concreta, busca
caracterizar, mediante una escalera comparativa incremental de cinco tratamientos,
qué nivel de sofisticación (estimación distribuida, dinámica poblacional y política
de revisión) se requiere para formar coaliciones de cardinalidad mínima de forma
distribuida. Cada tratamiento
añade un mecanismo sobre el anterior y se evalúa qué problema resuelve y cuál
deja abierto: T1 (heurística voraz local), T2 (consenso dinámico de media más
regla voraz), T3 (replicador distribuido modificado), T4 (dinámicas de Smith con
tasa de revisión espacialmente modulada) y T5 (referencia centralizada
replanificada). El método propuesto en este TFM es el tratamiento~T4; el
contraste entre su tasa de revisión modulada y una tasa constante se trata como
ablación interna, sin constituir un tratamiento separado. T5 funciona como
referencia centralizada que acota empíricamente el coste de la descentralización,
no como un óptimo absoluto. La metodología se organiza en dieciséis subsecciones que cubren
el modelo del sistema, las dinámicas acopladas, el análisis de convergencia, el
diseño experimental, los criterios de aceptación y una extensión de capacidad
heterogénea (formulación).

\subsection{Modelo del sistema}

El modelo del sistema se compone de tres elementos: los robots, descritos por una
cinemática de uniciclo; las tareas, caracterizadas por su posición y su requisito de
cardinalidad mínima; y el grafo de comunicación, que determina la información local
disponible para cada robot. Los tres se formalizan a continuación y fijan la notación
empleada en el resto del capítulo.

\subsubsection{Robots}

Cada AGV sigue cinemática uniciclo en tiempo discreto con período $T_s = 0{,}1\,\text{s}$:
\begin{align}
  x_i(t+1) &= x_i(t) + T_s \, v_i(t)\cos\theta_i(t), \label{eq:unicycle_x}\\
  y_i(t+1) &= y_i(t) + T_s \, v_i(t)\sin\theta_i(t), \label{eq:unicycle_y}\\
  \theta_i(t+1) &= \theta_i(t) + T_s \, \omega_i(t), \label{eq:unicycle_theta}
\end{align}
con entrada de control $(v_i, \omega_i)$ sujeta a $|v_i| \leq v_{\max}$ y
$|\omega_i| \leq \omega_{\max}$. El estado del robot $i$ es
$\mathbf{q}_i(t) = (x_i, y_i, \theta_i) \in \RR^3$. La Figura~\ref{fig:robot}
ilustra este modelo de tracción diferencial, sus entradas $v_i$ y $\omega_i$ y el
punto adelantado $\mathbf{h}_i$ empleado por la navegación reactiva.

\begin{figure}[htbp]
\centering
\begin{tikzpicture}[>={Stealth[length=5pt]}, font=\small, x=1cm, y=1cm]
  %-- Marco del mundo
  \draw[->, VIUGray!80] (0,0) -- (1.4,0) node[right]{$x$};
  \draw[->, VIUGray!80] (0,0) -- (0,1.3) node[above]{$y$};
  \node[VIUGray!80, below left=-1pt] at (0,0) {$O$};

  \def\cx{4.7}\def\cy{2.0}\def\ang{30}\def\rb{0.85}

  %-- Linea de referencia y arco del angulo de orientacion
  \draw[densely dotted, VIUGray] (\cx,\cy) -- (\cx+2.7,\cy);
  \draw[->, VIUGray] (\cx+1.55,\cy) arc (0:\ang:1.55);
  \node[VIUGray] at (\cx+1.82,\cy+0.42) {$\theta_i$};

  \begin{scope}[shift={(\cx,\cy)}, rotate=\ang]
    %-- Ruedas (a lo largo del eje, perpendiculares al avance)
    \fill[VIUDark] (-0.14,  \rb-0.02) rectangle (0.14,  \rb+0.48);
    \fill[VIUDark] (-0.14, -\rb-0.48) rectangle (0.14, -\rb+0.02);
    %-- Eje entre ruedas
    \draw[VIUDark!75, thick] (0,-\rb) -- (0,\rb);
    %-- Cuerpo del robot
    \draw[very thick, draw=VIUOrange, fill=VIUOrange!12] (0,0) circle (\rb);
    %-- Velocidad angular (rotacion)
    \draw[->, thick, VIUOrange] ([shift=(55:0.5)]0,0) arc (55:160:0.5);
    \node[VIUOrange] at ([shift=(108:0.78)]0,0) {$\omega_i$};
    %-- Velocidad lineal (avance)
    \draw[->, very thick, VIUDeep] (0,0) -- (\rb+1.0,0);
    \node[VIUDeep, anchor=south] at (\rb+0.55,0.05) {$v_i$};
    %-- Punto adelantado h_i a distancia ell del centro
    \draw[densely dashed, VIULinkBlue, thick] (\rb,0) -- (\rb+1.95,0);
    \fill[VIULinkBlue] (\rb+1.95,0) circle (2pt);
    \node[VIULinkBlue, anchor=south] at (\rb+1.95,0.07) {$\mathbf{h}_i$};
    %-- Medida de ell (del centro al punto adelantado)
    \draw[<->, VIUGray] (0,-1.45) -- (\rb+1.95,-1.45);
    \node[VIUGray, fill=white, inner sep=1pt] at (1.4,-1.45) {$\ell$};
  \end{scope}

  %-- Centro y etiqueta de pose
  \fill[VIUDark] (\cx,\cy) circle (1.8pt);
  \node[VIUDark, anchor=north east] at (\cx-0.05,\cy-0.08) {$(x_i,y_i)$};
\end{tikzpicture}
\caption[Robot móvil de tracción diferencial (uniciclo)]%
{Modelo de robot móvil de tracción diferencial (uniciclo). El estado es la pose
$\mathbf{q}_i=(x_i,y_i,\theta_i)$; las entradas de control son la velocidad lineal
$v_i$ y la velocidad angular $\omega_i$. El punto adelantado $\mathbf{h}_i$, situado
a una distancia $\ell$ del centro en la dirección de avance, es el que regula la ley
de navegación reactiva (ec.~\ref{eq:lookahead}).}
\label{fig:robot}
\end{figure}


\subsubsection{Tareas}

Cada tarea $k \in \{1, \ldots, K\}$ queda definida por su posición de recogida
$\mathbf{l}_k \in \RR^2$, su destino de entrega $\mathbf{d}_k \in \RR^2$,
su recompensa $r_k > 0$ y su requisito mínimo de cardinalidad $n_k \in \NN$.
La tarea $k = 0$ representa el estado de inactividad, con $r_0 = 0$ y $n_0 = 0$. El conjunto de estrategias disponibles para cada robot es
$\mathcal{S} = \{0, 1, \ldots, K\}$.

La cardinalidad mínima $n_k$ es un caso particular de una restricción más general de
capacidad acumulada: si cada robot $i$ posee una capacidad $c_i > 0$ y la carga $k$
demanda una capacidad $M_k$, la factibilidad se expresa como
$\sum_{i\in\mathcal{C}_k} c_i \geq M_k$. El requisito de cardinalidad se recupera con
flota homogénea ($c_i = c$, $M_k = n_k\,c$). El núcleo formal y experimental de este
TFM supone flota homogénea; la extensión heterogénea se formula en
\S\,\ref{subsec:heterogenea}.

\subsubsection{Grafo de comunicación}

El grafo de comunicación $G(t) = (\mathcal{V}, \mathcal{E}(t))$ es dinámico:
dos robots $i$ y $j$ son vecinos si $\norm{\mathbf{q}_i^{xy}(t) - \mathbf{q}_j^{xy}(t)}
\leq R$, donde $\mathbf{q}_i^{xy} = (x_i, y_i)$ y $R > 0$ es el rango de
comunicación. El conjunto de vecinos del robot $i$ en el instante $t$ es
$\mathcal{N}_i(t) = \{j \neq i : \norm{\mathbf{q}_i^{xy}(t) - \mathbf{q}_j^{xy}(t)} \leq R\}$.
El grado máximo del grafo es $d_{\max}(t) = \max_i |\mathcal{N}_i(t)|$.

\subsection{Función de payoff multiplicativa}

La función de payoff del robot $i$ para la tarea $k \geq 1$ tiene la forma
multiplicativa:
\begin{equation}
  f_{ik}(\hat{x}_i, \mathbf{q}_i) = r_k \cdot \Phi\!\left(N\hat{x}_{ik},\, n_k\right)
    \cdot \Psi\!\left(\norm{\mathbf{q}_i^{xy} - \mathbf{l}_k}\right),
  \label{eq:payoff}
\end{equation}
donde los dos factores que componen el payoff son:

\begin{equation}
  \Phi(z, n) = \frac{1}{1 + \exp\!\left(\beta(z - n)\right)},
  \quad \beta > 0,
  \label{eq:sigmoid}
\end{equation}
\begin{equation}
  \Psi(d) = \exp\!\left(-\frac{d^2}{2\lambda^2}\right),
  \quad \lambda > 0.
  \label{eq:gaussian}
\end{equation}

El factor $\Phi(N\hat{x}_{ik}, n_k)$ es la \emph{demanda sigmoide de coalición}.
Es monótonamente decreciente en $z = N\hat{x}_{ik}$ (el número estimado de
robots en la tarea $k$): cuando la tarea tiene pocos robots respecto a $n_k$,
$\Phi \approx 1$ y el payoff es alto (la tarea \emph{recluta} ag\-entes);
cuando tiene suficientes, $\Phi \approx 0$ y el payoff cae (la tarea está
\emph{saturada}). Esta monotonicidad decreciente es la que preserva la estructura
de juego de congestión ponderado y distingue este diseño del planteamiento
de cotas superiores de \citet{martinezpiazuelo2022automatica}.

El factor $\Psi(\cdot)$ es el descuento espacial gaussiano: robots lejanos de la
tarea perciben un payoff reducido, lo que incorpora la dinámica espacial al
juego de asignación. Para el estado de inactividad se asigna un payoff base
positivo pequeño:
\begin{equation}
  f_{i0} = r_0, \quad r_0 > 0 \text{ pequeño (por defecto } r_0 = 0{,}05),
  \label{eq:payoff_idle}
\end{equation}
lo que evita que la dinámica de Smith se congele por ausencia de diferencias de
payoff cuando todas las tareas están saturadas ($\Phi \approx 0$ para todo $k \geq 1$).
En el juego base con $\Psi \equiv 1$, tareas homogéneas y $N = \sum_k n_k$, la
condición $r_0 < r/2$ hace que $x_0^* = 0$ sea compatible con el equilibrio. En el
sistema espacial completo, donde $\Psi(d) < 1$, la inactividad puede ser localmente
preferible para robots lejanos o excedentes; en ese caso la inactividad se interpreta
como estrategia de reserva para robots no comprometidos. La caracterización del
equilibrio fuera del caso base se evalúa empíricamente.

\begin{figure}[htbp]
\centering
\begin{tikzpicture}[x=1cm,y=1cm]

%===== Panel (a): Phi(z,n) ============================
\begin{scope}[xshift=0cm]
  % Axes
  \draw[-{Stealth[length=5pt]}, VIUDark!70] (-0.2,0) -- (6.4,0)
    node[right, font=\small] {$z=N\hat{x}_{ik}$};
  \draw[-{Stealth[length=5pt]}, VIUDark!70] (0,-0.2) -- (0,3.4)
    node[above, font=\small] {$\Phi(z,n_k)$};

  % Sigmoid curve: Phi=1/(1+exp(2*(z-3))), scaled by 3
  \draw[VIUOrange, very thick, smooth]
    plot[domain=0:6, samples=60, variable=\z] (\z, {3/(1+exp(2*(\z-3)))});

  % Dashed reference lines at z=n=3
  \draw[VIUGray, dashed, thin] (3,0) node[below, font=\tiny] {$n_k$} -- (3,1.5);
  \draw[VIUGray, dashed, thin] (0,1.5) node[left, font=\tiny] {$\tfrac{1}{2}$} -- (3,1.5);
  \fill[VIUGray] (3,1.5) circle(1.5pt);

  % Tick at z=0: Phi=1 (approximately)
  \draw[VIUGray, thin] (-0.05,3) -- (0.05,3);
  \node[left, font=\tiny] at (0,3) {$1$};
  \draw[VIUGray, thin] (-0.05,0) -- (0.05,0);
  \node[left, font=\tiny] at (0,0) {$0$};

  % Zone annotations
  \draw[{Stealth[length=3pt]}-{Stealth[length=3pt]}, VIUDark!60, thin] (0.1,2.7) -- (2.9,2.7);
  \node[font=\tiny, VIUDark, align=center] at (1.5,3.15)
    {Zona de\\reclutamiento};

  \draw[{Stealth[length=3pt]}-{Stealth[length=3pt]}, VIUDark!60, thin] (3.1,0.25) -- (5.9,0.25);
  \node[font=\tiny, VIUDark, align=center] at (4.5,0.65)
    {Zona de\\saturación};

  % beta label on curve
  \node[font=\tiny, VIUOrange, anchor=west] at (3.3,2.3) {$\beta=2$};

  % Panel label
  \node[font=\small\bfseries, VIUDark] at (3.0,-0.85)
    {(a) Demanda sigmoide $\Phi(z,n_k)$};
\end{scope}

%===== Panel (b): Psi(d) =================================
\begin{scope}[xshift=8.0cm]
  % Axes
  \draw[-{Stealth[length=5pt]}, VIUDark!70] (-0.2,0) -- (6.4,0)
    node[right, font=\small] {$d=\|\mathbf{q}_i^{xy}-\mathbf{l}_k\|$};
  \draw[-{Stealth[length=5pt]}, VIUDark!70] (0,-0.2) -- (0,3.4)
    node[above, font=\small] {$\Psi(d)$};

  % Gaussian: Psi=exp(-d^2/50), x-axis 0..10m plotted as 0..6
  % scale: x_plot = d * 6/10; so d = x_plot * 10/6
  % Psi = exp(-(x_plot*10/6)^2/50) = exp(-x_plot^2/18)
  \draw[VIUDeep, very thick, smooth]
    plot[domain=0:6, samples=60, variable=\d] (\d, {3*exp(-\d*\d/18)});

  % Tick marks
  \draw[VIUGray, thin] (-0.05,3) -- (0.05,3);
  \node[left, font=\tiny] at (0,3) {$1$};
  \draw[VIUGray, thin] (-0.05,0) -- (0.05,0);
  \node[left, font=\tiny] at (0,0) {$0$};

  % lambda marker
  \draw[VIUGray, dashed, thin] (3,0) -- (3,{3*exp(-3*3/18)});
  \draw[VIUGray, dashed, thin] (0,{3*exp(-9/18)}) -- (3,{3*exp(-9/18)});
  \node[below, font=\tiny] at (3,0) {$\lambda$};
  \node[left, font=\tiny] at (0,{3*exp(-0.5)}) {$e^{-1/2}$};
  \fill[VIUGray] (3,{3*exp(-9/18)}) circle(1.5pt);

  % Annotation
  \node[font=\tiny, VIUDeep, anchor=west] at (3.3,2.4) {$\lambda=5\,\mathrm{m}$};

  % Panel label
  \node[font=\small\bfseries, VIUDark] at (3.0,-0.85)
    {(b) Descuento espacial $\Psi(d)$};

  % x-axis labels
  \node[below, font=\tiny] at (0,0) {$0$};
  \node[below, font=\tiny] at (6,0) {$10\,\mathrm{m}$};
\end{scope}

\end{tikzpicture}
\caption[Componentes de la función de payoff]%
{Componentes de la función de payoff \eqref{eq:payoff}. \textbf{(a)}~La
demanda sigmoide $\Phi(z,n_k)$ es monótonamente decreciente: cuando el número
estimado $z=N\hat{x}_{ik}$ está por debajo de $n_k$, el payoff es alto (zona de
reclutamiento); cuando supera $n_k$, cae a cero (zona de saturación). \textbf{(b)}~El
descuento gaussiano $\Psi(d)$ reduce el payoff de robots lejanos de la tarea, integrando
la dinámica espacial en el juego de asignación.}
\label{fig:payoff}
\end{figure}


\subsection{Dinámica~1: consenso dinámico de media (DAC)}

Cuando la estrategia $s_i(t)$ cambia, la fracción de robots en cada tarea cambia
también; el estimador debe seguir esa señal variable sin reinicializaciones
abruptas que introduzcan discontinuidades en el payoff.
Se define la señal local $y_{ik}(t) = \mathbf{1}_{[s_i(t)=k]}$ y se emplea el
\emph{consenso dinámico de media} (DAC) \citep{kia2019tutorial}:
\begin{equation}
  \hat{x}_{ik}(t+1) = \hat{x}_{ik}(t)
    + \varepsilon \!\sum_{j \in \mathcal{N}_i(t)}\!\bigl[\hat{x}_{jk}(t) - \hat{x}_{ik}(t)\bigr]
    + \bigl(y_{ik}(t+1) - y_{ik}(t)\bigr),
  \quad \forall k \in \mathcal{S},
  \label{eq:consenso}
\end{equation}
con ganancia $\varepsilon \in (0, 1/\Delta_{\max})$, donde
$\Delta_{\max} = \max_t d_{\max}(t)$; cuando $d_{\max}$ no se conoce a priori se
adopta la cota conservadora $\Delta_{\max} = N-1$, válida para cualquier grafo sobre
$N$ nodos. Una $\varepsilon$ conservadora ralentiza el consenso y afecta por igual a
T2 y T3 (que dependen del DAC); este efecto se tiene en cuenta al interpretar la
comparación con T4. El tercer término
$\Delta y_{ik}(t) = y_{ik}(t+1) - y_{ik}(t)$ compensa instantáneamente el salto
en la señal local cuando el robot reasigna su estrategia, sin necesidad de
reinicializar el estado del estimador.

El término $\Delta y_{ik}$ requiere conocer $s_i(t+1)$ antes de actualizar
$\hat{x}_{ik}$; por ello la dinámica de Smith (Bloque~2) precede a la actualización
del estimador (Bloque~3) en el bucle de control (véase Algoritmo~\ref{alg:loop}).

El payoff se evalúa con el argumento clippeado al intervalo válido:
\begin{equation}
  f_{ik}(\hat{x}_i, \mathbf{q}_i)
    = r_k \cdot \Phi\!\bigl(N\,\mathrm{clip}(\hat{x}_{ik},0,1),\,n_k\bigr)
      \cdot \Psi\!\bigl(\norm{\mathbf{q}_i^{xy} - \mathbf{l}_k}\bigr),
  \label{eq:payoff_clip}
\end{equation}
donde $\mathrm{clip}(a,0,1) = \min(1,\max(0,a))$.
El estado interno $\hat{x}_{ik}$ puede tomar valores fuera de $[0,1]$ de forma
transitoria (efecto normal del DAC durante reasignaciones masivas) sin afectar la
estabilidad del bucle, porque el payoff resultante se mantiene acotado en $(0, r_k]$.

Se distinguen tres medidas de ocupación de la tarea $k$: la \emph{ocupación
asignada pura} $a_k(t) = \sum_i \mathbf{1}_{[s_i(t)=k]}$; la \emph{ocupación física
presente} $\mathrm{count}_k(t) = \sum_i \mathbf{1}_{[s_i(t)=k]}\,
\mathbf{1}_{[\norm{\mathbf{q}_i^{xy}(t)-\mathbf{l}_k}<r_{\mathrm{det}}]}$, que coincide
con $|\mathcal{C}_k(t)|$; y la \emph{estimación local} $\hat{x}_{ik}(t)$ obtenida por
DAC. La estimación $\hat{x}_{ik}$ alimenta el payoff; $a_k$ mide la asignación lógica;
y $\mathrm{count}_k = |\mathcal{C}_k|$ activa la fase de transporte y determina la
factibilidad física $|\mathcal{C}_k| \geq n_k$.

\begin{observacion}[Error de seguimiento del DAC y sensibilidad al payoff]
\label{obs:init_obs}
Para un grafo estacionario conectado con $N$ agentes, el DAC \eqref{eq:consenso}
mantiene el error de seguimiento acotado por $\|\hat{x}_{ik}(t) - x_k(t)\| = O(K/N)$
bajo cambios de estrategia individuales y recupera la media exacta en tiempo
$O(1/(\varepsilon\lambda_2))$ entre eventos \citep{kia2019tutorial}.
La inicialización uniforme $\hat{x}_{ik}(0) = 1/(K+1)$ acelera la convergencia inicial;
la indicador $\hat{x}_{ik}(0) = \mathbf{1}_{[s_i(0)=k]}$ es igualmente válida
pues el DAC absorbe el salto inicial como si fuera un evento de reasignación.
En los experimentos se comparan ambas inicializaciones en el Escenario~A.
\end{observacion}
En grafos dinámicos con $R$ finito, el error de consenso $e_{ik}(t) = \hat{x}_{ik}(t)
- x_k(t)$ introduce un sesgo acotado en los payoffs, cuyo efecto sobre el equilibrio
se cuantifica en los experimentos del Escenario~D.

\subsection{Dinámica~2: Smith distribuidas en tiempo discreto}
\label{subsec:smith}

El vector de preferencias $\mathbf{p}_i(t)$ evoluciona según las dinámicas de
Smith en tiempo discreto:
\begin{equation}
  \tilde{p}_{ik}(t+1) = p_{ik}(t) + T_s \mu_i(t) \left[
    \sum_{j \in \mathcal{S}} p_{ij}(t)\pos{f_{ik} - f_{ij}}
    - p_{ik}(t) \sum_{j \in \mathcal{S}} \pos{f_{ij} - f_{ik}}
  \right],
  \label{eq:smith_update}
\end{equation}
\begin{equation}
  \mathbf{p}_i(t+1) = \Pi_{\simplex}\!\left(\tilde{\mathbf{p}}_i(t+1)\right),
  \label{eq:proyeccion}
\end{equation}
donde $\pos{y} = \max(y, 0)$, $\Pi_{\simplex}$ es la proyección ortogonal sobre
el símplex $\simplex^K$ (calculada por el algoritmo estándar de clasificación
y umbral de complejidad $O(K\log K)$), y $\mu_i(t) > 0$ es la tasa de revisión del robot $i$.
La proyección $\Pi_{\simplex}$ se implementa como salvaguarda numérica frente al
error de punto flotante; la cota de paso $\beta_{\max}$ ya garantiza la invarianza
del símplex en la dinámica ideal.
El primer término de \eqref{eq:smith_update} representa el flujo de entrada desde
estrategias con menor payoff; el segundo, el flujo de salida hacia estrategias con
mayor payoff. Esta comparación bilateral elimina la necesidad de conocer la media
global del payoff, a diferencia del replicador distribuido.

\paragraph{Regla de histéresis para la selección de estrategia.}
La estrategia efectiva del robot se actualiza con un umbral de histéresis
$\rho_s > 0$ (valor nominal $\rho_s = 0{,}05$, Tabla~\ref{tab:params}):
\begin{equation}
  s_i(t+1) = \begin{cases}
    \displaystyle\argmax_{k \in \mathcal{S}}\, p_{ik}(t+1) &
      \text{si } p_{ik^*}(t+1) - p_{i s_i(t)}(t+1) > \rho_s, \\
    s_i(t) & \text{en otro caso,}
  \end{cases}
  \label{eq:hysteresis}
\end{equation}
donde $k^* = \argmax_k p_{ik}(t+1)$. La condición exige que la estrategia dominante
aventaje a la actual en al menos $\rho_s$ antes de ejecutar el cambio, lo que reduce
el \emph{chattering} sin alterar los puntos de equilibrio de la dinámica.

\begin{lema}[Factibilidad de la partición pura]
\label{lema:factibilidad}
Bajo la regla \eqref{eq:hysteresis} con $\rho_s \geq 0$, los conjuntos $\{s_i(t)\}$
forman en todo instante una partición pura de los $N$ robots sobre $\mathcal{S}$
(incluyendo la tarea de inactividad $k=0$). En particular, $\sum_{k=0}^K |\mathcal{C}_k(t)| = N$
en todo $t$: la cardinalidad total de las coaliciones es siempre exactamente $N$.
\end{lema}
\begin{proof}
La regla \eqref{eq:hysteresis} es determinista y asigna cada robot a exactamente un
$s_i(t) \in \mathcal{S}$ en cada paso, con empate resuelto manteniendo $s_i(t)$.
La partición se preserva por inducción trivial. La igualdad de la suma se sigue
de que los $N$ robots son disjuntos.
\end{proof}

\subsection{Tasa de revisión espacialmente modulada}

La tasa de revisión de cada robot depende de su distancia a la posición de la
tarea que actualmente persigue:
\begin{equation}
  \mu_i(t) = \mu_{\min} + (\mu_0 - \mu_{\min})
    \left(1 - \exp\!\left(-\frac{\norm{\mathbf{q}_i^{xy}(t) -
    \mathbf{l}_{s_i(t)}}^2}{2\sigma_{\text{rev}}^2}\right)\right),
  \label{eq:revision_rate}
\end{equation}
con $\mu_0 > 0$, $\mu_{\min} = 0{,}02\,\mu_0$ y $\sigma_{\text{rev}} > 0$.
El piso $\mu_{\min}$ mantiene una tasa mínima de revisión incluso en robots muy
próximos al destino, evitando que el lazo de Smith quede congelado ante perturbaciones. Los robots lejanos tienen $\mu_i \approx \mu_0$:
pueden reconsiderar con tasa plena. Para robots en estado de inactividad
($s_i = 0$), se usa $\mu_i = \mu_0$.

La tasa modulada actúa como un término de histéresis suave que mitiga el
\emph{chattering} cerca de los umbrales de cardinalidad $n_k$: los robots comprometidos
con su tarea son más resistentes a la reasignación que los que aún navegan
hacia su destino.

\subsection{Dinámica~3: campo de gradiente reactivo}

La navegación regula prácticamente un \emph{punto adelantado} $\mathbf{h}_i$
situado a distancia $\ell$ por delante del centro del robot en su dirección de
avance:
\begin{equation}
  \mathbf{h}_i = \mathbf{q}_i^{xy} + \ell\,(\cos\theta_i,\ \sin\theta_i)^{\top},
  \qquad \ell > 0.
  \label{eq:lookahead}
\end{equation}
El movimiento se obtiene minimizando localmente un campo de potencial artificial
\citep{bullo2009distributed} definido sobre el punto adelantado:
\begin{equation}
  F_i(\mathbf{h}_i) = \frac{\alpha}{2}\norm{\mathbf{h}_i - \mathbf{l}_{s_i}^{\mathrm{eff}}}^2
    + \eta \sum_{j \in \mathcal{N}_i}
      \frac{1}{\norm{\mathbf{q}_i^{xy} - \mathbf{q}_j^{xy}}^2 + \varepsilon_r^2}
    + \delta \sum_{j \in \mathcal{C}_k \setminus \{i\}} \frac{(\norm{\mathbf{h}_i
    - \mathbf{h}_j} - d^*)^2}{2}
    \cdot \mathbf{1}_{[\norm{\mathbf{h}_i - \mathbf{l}_k} < r_{\text{form}}]},
  \label{eq:potential}
\end{equation}
donde el primer término crea un pozo de potencial con mínimo en el objetivo
efectivo $\mathbf{l}_{s_i}^{\mathrm{eff}}$ (ec.~\eqref{eq:target_switch}); el
segundo repele a los robots entre sí (evaluado sobre los centros físicos
$\mathbf{q}_i^{xy}$, para reflejar la geometría real de colisión); y el tercero
mantiene la geometría de coalición. El término atractivo y la conversión cinemática
operan sobre el punto adelantado $\mathbf{h}_i$; la repulsión inter-robot, en cambio,
se calcula sobre los centros físicos $\mathbf{q}_i^{xy}$. El parámetro $\varepsilon_r = 0{,}01\,\text{m}$
regulariza el denominador cuadrático y evita singularidades; la forma
$\norm{\cdot}^2 + \varepsilon_r^2$ es diferenciable en todo $\mathbb{R}^2$ y
proporciona un gradiente repulsivo acotado incluso para robots en contacto. Los
campos de potencial artificial presentan el problema clásico de mínimos locales
\citep{bullo2009distributed}; en los experimentos se mitiga mediante una
perturbación aleatoria cuando la velocidad del robot cae por debajo de
$0{,}01\,\text{m/s}$ durante más de 5 pasos consecutivos. El término de repulsión
no garantiza formalmente la ausencia de colisiones: la evitación es heurística y
depende de los valores de $\eta$ y $\varepsilon_r$. La distancia geométrica de
formación es $d^* = 2(r_{\text{load}} + r_{\text{robot}}) \sin(\pi / n_{s_i(t)})$,
donde $n_{s_i(t)} \equiv n_k$ con $k = s_i(t)$ es la cardinalidad mínima de la
tarea actualmente asignada al robot $i$.
Para $n_k = 1$ la fórmula da $d^* = 0$ (sin geometría de formación, el robot
se posiciona directamente sobre la carga); para $n_k = 2$, $d^* = 2(r_{\text{load}}
+ r_{\text{robot}})$ (los dos robots quedan en lados opuestos de la carga). La zona
de formación se activa cuando $\norm{\mathbf{h}_i - \mathbf{l}_k} < r_{\text{form}}$.
Análogamente a $\varepsilon_r$, el gradiente del término de formación emplea una
regularización $\varepsilon_f$ en el denominador del vector unitario,
$(\mathbf{h}_i-\mathbf{h}_j)/(\norm{\mathbf{h}_i-\mathbf{h}_j}+\varepsilon_f)$, lo que
evita la singularidad cuando $\norm{\mathbf{h}_i-\mathbf{h}_j}\to 0$.

Cuando el escenario incluye obstáculos estáticos $\mathcal{O} = \{O_1,\ldots,O_M\}$,
a los términos de \eqref{eq:potential} se añade un término repulsivo de obstáculos
\begin{equation}
  F_i^{\mathrm{obs}} = \eta_o \sum_{m=1}^{M}
    \frac{1}{\mathrm{dist}(\mathbf{q}_i^{xy}, O_m)^2 + \varepsilon_o^2},
  \label{eq:obstaculo}
\end{equation}
calculado sobre el centro físico $\mathbf{q}_i^{xy}$, donde $\varepsilon_o$ regulariza
el denominador. Es una heurística reactiva de simulación que no sustituye la
planificación global de rutas ni certifica la ausencia de colisiones.

El campo de gradiente se aplica sobre el punto adelantado, no sobre el centro, y se
convierte en comandos cinemáticos invirtiendo el jacobiano del punto adelantado:
\begin{equation}
  \begin{bmatrix} v_i \\ \omega_i \end{bmatrix}
    = \begin{bmatrix} \cos\theta_i & -\ell\sin\theta_i \\
                      \sin\theta_i & \phantom{-}\ell\cos\theta_i \end{bmatrix}^{-1}
      \mathbf{u}_h, \qquad \mathbf{u}_h = -\nabla_{\mathbf{h}_i} F_i,
  \label{eq:lookahead_cmd}
\end{equation}
seguida de saturación a $|v_i| \leq v_{\max}$ y $|\omega_i| \leq \omega_{\max}$. El
jacobiano tiene determinante $\ell$, por lo que es invertible para todo $\ell > 0$.
En la práctica, la conversión regula el punto adelantado $\mathbf{h}_i$ a una
vecindad del objetivo y evita aplicar un gradiente cartesiano directamente sobre el
centro geométrico del robot diferencial, con lo que respeta su naturaleza no holonómica.

\subsubsection{Fase de transporte y transición al destino}
\label{subsubsec:transporte}

Una vez que la coalición $\mathcal{C}_k$ cumple $|\mathcal{C}_k(t)| \geq n_k$
durante al menos $\tau_s = 3$ pasos consecutivos, el sistema activa el
\emph{modo de transporte}: el objetivo del campo de gradiente \eqref{eq:potential}
cambia de la posición de recogida $\mathbf{l}_k$ al destino de descarga $\mathbf{d}_k$.
Formalmente, el término de atracción de \eqref{eq:potential} emplea el objetivo efectivo:
\begin{equation}
  \mathbf{l}_{s_i}^{\mathrm{eff}}(t) =
  \begin{cases}
    \mathbf{l}_k & \text{si } |\mathcal{C}_k(t)| < n_k \;\text{(fase de reclutamiento)},\\
    \mathbf{d}_k & \text{si } |\mathcal{C}_k(t)| \geq n_k \;\text{(fase de transporte)}.
  \end{cases}
  \label{eq:target_switch}
\end{equation}

La transición de modo~\eqref{eq:target_switch} supone
$R > 2(r_{\text{load}} + r_{\text{robot}})$, condición que garantiza que los robots
capaces de rodear una misma carga forman un subgrafo completo y evalúan
$|\mathcal{C}_k| \geq n_k$ de forma sincronizada; por debajo de ese umbral la
decisión de modo puede fragmentarse, régimen caracterizado en el Escenario~D.

Además, durante la fase de transporte el factor espacial $\Psi(\cdot)$ se fija
en $1{,}0$ para todos los robots de $\mathcal{C}_k$. Sin esta corrección, la
distancia $\norm{\mathbf{q}_i^{xy} - \mathbf{l}_k}$ crecería a medida que la
coalición se aleja del punto de recogida hacia $\mathbf{d}_k$, lo que reduciría
artificialmente el payoff de la tarea y podría inducir defecciones de la
coalición durante el transporte.

Durante la fase de transporte, el centroide de la coalición $\mathcal{C}_k$ navega
hacia el destino $\mathbf{d}_k$ bajo el campo de potencial, con el término de
obstáculos $F_i^{\mathrm{obs}}$ \eqref{eq:obstaculo} activo sobre cada robot; los
miembros de la coalición mantienen la geometría de formación $d^*$ relativa al
centroide en movimiento. Así, la coalición desvía su trayectoria
colectiva ante obstáculos estáticos sin disolverse. El modelo permanece
cinemático: no se simulan fuerzas de contacto sobre la carga. La evitación a nivel de
coalición se evalúa empíricamente en los escenarios con obstáculos y cualitativamente en
CoppeliaSim (Escenario~F).

La transición es bidireccional: si durante el modo de transporte la coalición
pierde un miembro y satisface $|\mathcal{C}_k(t)| < n_k$ durante $\tau_s$ pasos
consecutivos (por ejemplo, por la retirada abrupta de un robot), el modo revierte
a reclutamiento, el objetivo efectivo regresa a $\mathbf{l}_k$ y el factor
$\Psi(\cdot)$ vuelve a calcularse dinámicamente. Esta histéresis evita
transiciones espurias debidas a fluctuaciones momentáneas en la coalición.

La tarea $k$ se considera completada cuando el centroide de $\mathcal{C}_k$ se
encuentra dentro de $r_{\mathrm{det}}$ del destino $\mathbf{d}_k$; en ese instante
se detiene el reloj de $\bar{T}$.

\begin{observacion}[Modelo cinemático de la fase de transporte]
El modelo de transporte es puramente cinemático: se propaga el movimiento
del \emph{centroide de la coalición} hacia $\mathbf{d}_k$, sin simular la
respuesta dinámica de la carga a las fuerzas de contacto. La geometría de
formación $d^*$ define las condiciones iniciales de posición relativa
que un controlador de manipulación aguas abajo (como el propuesto por
\citealt{ebel2024cooperative}) consumiría para calcular las fuerzas necesarias.
En la simulación Python se registra la trayectoria de la pose de la carga como
el centroide de $\mathcal{C}_k$ durante el modo de transporte; en la validación
CoppeliaSim (Escenario~F) se representa como un cuboide pasivo acoplado
cinemáticamente a ese centroide.
\end{observacion}

\subsection{Acoplamiento entre las tres dinámicas}
\label{subsec:acoplamiento}

Las tres dinámicas están bidireccionalmente acopladas en un único bucle de
control ejecutado cada $T_s = 0{,}1\,\text{s}$ (Fig.~\ref{fig:loop}). Este
acoplamiento forma un sistema retroalimentado en el que la posición afecta
al payoff, el payoff afecta a la estrategia, la estrategia afecta al movimiento,
y el movimiento cambia la posición y la topología de comunicación.

\begin{figure}[htbp]
\centering
\begin{tikzpicture}[
  font=\small,
  block/.style={
    draw=VIUDark!55, very thick, rounded corners=4pt,
    minimum width=5.4cm, minimum height=1.1cm,
    align=center, fill=VIULight, inner sep=4pt,
  },
  flow/.style={-{Latex[length=6pt,width=5pt]}, line width=1.1pt, draw=VIUOrange},
  ret/.style={-{Latex[length=6pt,width=5pt]}, line width=1.0pt, draw=VIUOrange!70, dashed},
  feed/.style={-{Latex[length=6pt,width=5pt]}, line width=1.1pt, draw=VIUDeep, dashed},
  lab/.style={font=\footnotesize, color=VIUDark, inner sep=1.5pt, fill=white, rounded corners=1.5pt},
  flab/.style={font=\footnotesize\itshape, color=VIUDeep, inner sep=1.5pt, fill=white, rounded corners=1.5pt},
]

% -------- Pipeline central (un paso temporal) --------
\node[block] (B1) at (0,0)
  {\textbf{1.\ Payoff}\\[-2pt]\footnotesize $f_{ik}=r_k\,\Phi(N\hat{x}_{ik},n_k)\,\Psi(d_{ik})$};
\node[block, below=12mm of B1] (B2)
  {\textbf{2.\ Dinámica de Smith}~(tasa $\mu_i$)\\[-2pt]\footnotesize $\mathbf{p}_i(t{+}1),\ s_i(t{+}1)$};
\node[block, below=12mm of B2] (B3)
  {\textbf{3.\ Consenso DAC}\\[-2pt]\footnotesize $\hat{x}_{ik}(t{+}1)$};
\node[block, below=12mm of B3] (B4)
  {\textbf{4.\ Navegación}~(gradiente + uniciclo)\\[-2pt]\footnotesize $\mathbf{q}_i(t{+}1)$};

\draw[flow] (B1) -- (B2) node[midway,lab]{$f_{ik}$};
\draw[flow] (B2) -- (B3) node[midway,lab]{$\Delta y_{ik}$};
\draw[flow] (B3) -- (B4) node[midway,lab]{$s_i(t{+}1)$};

% -------- Retorno de estimación: x̂ del DAC al payoff (izquierda) --------
\draw[ret] (B3.west) -- ++(-1.8,0) |- (B1.west)
  node[pos=0.25,left,lab]{$\hat{x}_{ik}(t{+}1)$};

% -------- Acoplamiento espacial: q de navegación a payoff (Ψ) y Smith (μ) --------
\coordinate (rt) at ($(B1.east)+(2.8,0)$);
\coordinate (rb) at ($(B4.east)+(2.8,0)$);
\draw[line width=1.1pt, draw=VIUDeep, dashed] (B4.east) -- (rb) -- (rt);
\draw[feed] (rt) -- (B1.east) node[midway,above,flab]{$\Psi(\mathbf{q}_i)$};
\draw[feed] ($(rt)!(B2.east)!(rb)$) -- (B2.east) node[midway,above,flab]{$\mu(\mathbf{q}_i)$};
\node[rotate=-90, font=\footnotesize\bfseries, color=VIUDeep, anchor=south]
  at ($(rb)!0.5!(rt)+(0.42,0)$) {acoplamiento espacial};

% -------- Etiqueta del paso temporal --------
\node[font=\footnotesize\itshape, color=VIUGray, anchor=west]
  at ($(B1.north west)+(0,0.55)$)
  {Paso $T_s = 0{,}1$\,s: decisión y movimiento en un único bucle};

\end{tikzpicture}
\caption[Bucle de control de las tres dinámicas acopladas]%
{Bucle de control de las tres dinámicas acopladas. Los cuatro bloques se ejecutan
secuencialmente dentro del mismo paso temporal $T_s = 0{,}1$\,s (flechas naranjas). La
estimación $\hat{x}_{ik}$ realimenta el payoff (retorno discontinuo izquierdo) y la
pose $\mathbf{q}_i$ modula el descuento espacial $\Psi$ del payoff y la tasa de
revisión $\mu_i$ de Smith (acoplamiento espacial, derecha): la decisión de asignación
y el movimiento quedan acoplados en un único bucle.}
\label{fig:loop}
\end{figure}


\begin{table}[htbp]
\centering
\caption[Acoplamiento entre dinámicas]{Acoplamiento bidireccional entre las tres dinámicas.}
\label{tab:acoplamiento}
\begin{tabularx}{\textwidth}{>{\bfseries}p{0.27\textwidth}Y}
\toprule
Interacción & Descripción \\
\midrule
$\mathbf{p}_i(t) \to f_{ik} \to s_i(t)$ &
  La estrategia determina el destino $\mathbf{l}_{s_i}$ y modifica el campo de
  potencial $F_i$. \\
$\mathbf{q}_i^{xy}(t) \to \Psi \to f_{ik}$ &
  La posición determina el descuento espacial del payoff de cada tarea. \\
$\mathbf{q}_i^{xy}(t) \to G(t) \to \hat{x}_{ik}$ &
  El movimiento cambia la topología del grafo y la velocidad de consenso. \\
$s_i(t{+}1) \to \Delta y_{ik} \to \hat{x}_{ik}$ &
  El cambio de estrategia inyecta instantáneamente $\Delta y_{ik}$ en el DAC,
  compensando el salto sin reinicializar el estimador. \\
$\mu_i(t) \gets \mathbf{q}_i, \mathbf{l}_{s_i}$ &
  La tasa de revisión depende de la distancia al destino actual. \\
\bottomrule
\end{tabularx}
\end{table}

\subsection{Pseudocódigo del bucle de control}
\label{subsec:pseudocodigo}

El Algoritmo~\ref{alg:loop} formaliza el bucle de control ejecutado por cada robot
$i$ en cada paso temporal. Los pasos de los tres bloques dinámicos se ejecutan
secuencialmente dentro del mismo instante $t$; las comunicaciones con los vecinos
$\mathcal{N}_i(t)$ se realizan al inicio del paso y se consideran atómicas.

\begin{algorithm}[htbp]
\caption{Bucle de control distribuido, robot $i$, instante $t$}
\label{alg:loop}
\begin{algorithmic}[1]
\Require $\mathbf{q}_i(t)$, $\mathbf{p}_i(t)$, $\hat{\mathbf{x}}_i(t)$,
         $s_i(t)$, $\{\mathbf{l}_k, \mathbf{d}_k, r_k, n_k\}$,
         $\mathcal{C}_k(t)$, modo$_k(t)$, contadores $\mathrm{cnt}^{+}_k, \mathrm{cnt}^{-}_k$
\Ensure $\mathbf{q}_i(t{+}1)$, $\mathbf{p}_i(t{+}1)$, $\hat{\mathbf{x}}_i(t{+}1)$

\State \textbf{Bloque~1: Cálculo de payoffs} \eqref{eq:payoff_clip}
\For{cada $k \geq 1$}
  \If{modo$_k(t) = $ transporte}
    \State $\mathbf{l}^{\mathrm{eff}}_k \gets \mathbf{d}_k$;\enspace
           $\Psi_k \gets 1{,}0$
  \Else
    \State $\mathbf{l}^{\mathrm{eff}}_k \gets \mathbf{l}_k$;\enspace
           $\Psi_k \gets \exp\!\bigl(-\norm{\mathbf{q}_i^{xy}{-}\mathbf{l}_k}^2/(2\lambda^2)\bigr)$
  \EndIf
  \State $f_{ik} \gets r_k \cdot \Phi\!\left(N\,\mathrm{clip}(\hat{x}_{ik}(t),0,1),\,n_k\right)\cdot\Psi_k$
\EndFor
\State $f_{i0} \gets r_0$

\State \textbf{Bloque~2: Dinámica de Smith} \eqref{eq:smith_update}--\eqref{eq:proyeccion}
\If{$s_i = 0$} \Comment{robot inactivo: sin objetivo espacial, tasa base}
  \State $\mu_i \gets \mu_0$
\Else
  \State $\mu_i \gets \mu_{\min} + (\mu_0{-}\mu_{\min})\!\left(1 - \exp\!\bigl(-\norm{\mathbf{q}_i^{xy}{-}\mathbf{l}^{\mathrm{eff}}_{s_i}}^2\!/(2\sigma_{\mathrm{rev}}^2)\bigr)\right)$
\EndIf
\For{cada $k \in \mathcal{S}$}
  \State $\tilde{p}_{ik} \gets p_{ik} + T_s\mu_i\!\left[
    \sum_{j}p_{ij}\pos{f_{ik}{-}f_{ij}}
    - p_{ik}\!\sum_{j}\pos{f_{ij}{-}f_{ik}}\right]$
\EndFor
\State $\mathbf{p}_i(t{+}1) \gets \Pi_{\simplex}(\tilde{\mathbf{p}}_i)$;\enspace
       $k^* \gets \argmax_k\, p_{ik}(t{+}1)$
\If{$p_{ik^*}(t{+}1) - p_{i,s_i(t)}(t{+}1) > \rho_s$}
  \State $s_i(t{+}1) \gets k^*$ \Comment{regla de histéresis~\eqref{eq:hysteresis}}
\Else
  \State $s_i(t{+}1) \gets s_i(t)$ \Comment{persiste bajo $\rho_s$ (incluye empates)}
\EndIf

\State \textbf{Bloque~3: Consenso dinámico de media (DAC)} \eqref{eq:consenso}
\For{cada $k \in \mathcal{S}$}
  \State $\hat{x}_{ik}(t{+}1) \gets \hat{x}_{ik}(t)
    + \varepsilon\!\sum_{j\in\mathcal{N}_i(t)}\!\bigl[\hat{x}_{jk}(t) - \hat{x}_{ik}(t)\bigr]
    + \bigl(y_{ik}(t{+}1) - y_{ik}(t)\bigr)$
    \Comment{$y_{ik}(t{+}1)=\mathbf{1}_{[s_i(t{+}1)=k]}$ proviene del Bloque~2}
\EndFor

\State \textbf{Bloque~4: Gradiente reactivo y cinemática} \eqref{eq:potential}--\eqref{eq:lookahead_cmd}
\State $\mathbf{h}_i \gets \mathbf{q}_i^{xy} + \ell\,(\cos\theta_i,\sin\theta_i)^{\top}$;\enspace
       $\mathbf{u}_h \gets -\nabla_{\mathbf{h}_i} F_i(\mathbf{h}_i;\,\mathbf{l}^{\mathrm{eff}}_{s_i})$
\State $[\,v_i;\ \omega_i\,]^{\top} \gets J(\theta_i,\ell)^{-1}\,\mathbf{u}_h$, saturado a $(v_{\max},\omega_{\max})$
       \Comment{$J(\theta_i,\ell)$: jacobiano del punto adelantado, ec.~\eqref{eq:lookahead_cmd}}
\State $\mathbf{q}_i(t{+}1) \gets $ unicycle\_step$(\mathbf{q}_i(t), v_i, \omega_i)$

\State \textbf{Actualización del grafo y modo de tarea}
\State $\mathcal{N}_i(t{+}1) \gets \{j\neq i : \norm{\mathbf{q}_i^{xy}(t{+}1){-}\mathbf{q}_j^{xy}(t{+}1)}\leq R\}$
\For{cada tarea $k$ con $i \in \mathcal{C}_k$}
  \If{$|\mathcal{C}_k(t{+}1)|\geq n_k$}
    \State $\mathrm{cnt}^{+}_k \gets \mathrm{cnt}^{+}_k + 1$;\enspace $\mathrm{cnt}^{-}_k \gets 0$
  \Else
    \State $\mathrm{cnt}^{-}_k \gets \mathrm{cnt}^{-}_k + 1$;\enspace $\mathrm{cnt}^{+}_k \gets 0$
    \Comment{reinicio al caer bajo $n_k$}
  \EndIf
  \If{modo$_k(t) = $ reclutamiento \textbf{y} $\mathrm{cnt}^{+}_k \geq \tau_s$}
    \State modo$_k(t{+}1) \gets $ transporte \Comment{ec.~\eqref{eq:target_switch}}
  \ElsIf{modo$_k(t) = $ transporte \textbf{y} $\mathrm{cnt}^{-}_k \geq \tau_s$}
    \State modo$_k(t{+}1) \gets $ reclutamiento \Comment{transición inversa}
  \EndIf
\EndFor
\end{algorithmic}
\end{algorithm}

\begin{observacion}[Error de seguimiento del DAC tras cambios simultáneos]
Cuando varios robots cambian de estrategia en el mismo instante $t$, los términos
$\Delta y_{ik}$ correspondientes inyectan el salto instantáneo en el estimador
sin reinicializarlo. Sin embargo, si $m$ robots reasignan simultáneamente, el DAC
acumula un error de seguimiento inicial acotado por $O(m/N)$; el transitorio decae
como $(1-\varepsilon\lambda_2)^t$ y se extingue en $O(1/(\varepsilon\lambda_2))$ pasos.
El efecto acumulado de estos transitorios queda capturado por la métrica
$\Delta_{\mathrm{chg}}$ en el Escenario~A.
\end{observacion}

\subsection{Análisis de convergencia}

El análisis formal se presenta en dos pasos: un teorema de convergencia para el
\emph{caso base ideal} y una proposición que acota la perturbación introducida por
el consenso local con $R$ finito. El análisis del caso general (grafo irregular,
tareas heterogéneas, cardinalidades distintas) queda como extensión doctoral.

\begin{teorema}[Convergencia, caso base ideal]
\label{prop:convergencia}
Considérese $N$ robots con rango de comunicación $R = \infty$ (sin error de
consenso), $K$ tareas homogéneas con $n_k = n$ para todo $k$ y $r_k = r$ para
todo $k$, $N = Kn$ (balance exacto), factor espacial $\Psi \equiv 1$ y sin
transición de modo. Bajo las dinámicas de Smith \eqref{eq:smith_update} con payoff
$f_{ik}(x) = r \cdot \Phi(Nx_k, n)$ (comunicación global perfecta, $\hat{x}_{ik} = x_k$),
donde $\Phi$ es estrictamente decreciente en $z = Nx_k$, el único
equilibrio Nash interior es:
\[
  x_k^* = \frac{n}{N} \quad \forall k \in \{1, \ldots, K\},
  \quad x_0^* = 0.
\]
Este equilibrio es localmente asintóticamente estable bajo la condición
\textup{(C1)}: $\beta \in (\beta_{\min}, \beta_{\max})$, donde los límites
están dados explícitamente por:
\begin{align}
  \beta_{\min} &= \ln\!\left(\frac{N+1}{N-1}\right),
    \label{eq:betamin}\\
  \beta_{\max} &= \frac{8}{T_s \cdot \mu_0 \cdot r \cdot N}.
    \label{eq:betamax}
\end{align}
El umbral $\beta_{\min}$ se introduce como criterio de diseño, no como condición
necesaria: garantiza una separación de payoff equivalente a un robot en la ocupación
normalizada ($1/N$) y previene gradientes desvanecientes en la zona de reclutamiento
de la sigmoide. La derivación siguiente formaliza ese criterio.
$\beta_{\min} \approx 2/N$ para $N$ grande (primer orden de Taylor).
$\beta_{\min}$ es una \emph{condición de separación de payoff}: garantiza que la diferencia $\Phi(n{-}1,n) - \Phi(n{+}1,n) > 1/N$ sea suficiente para distinguir el estado subtamaño ($z = n{-}1$)
del sobretamaño ($z = n{+}1$). La derivación es directa: con
$\Phi(z,n) = 1/(1+e^{\beta(z-n)})$ se tiene
$\Phi(n{-}1,n) - \Phi(n{+}1,n) = (e^\beta - 1)/(e^\beta + 1)$;
la condición $(e^\beta-1)/(e^\beta+1) > 1/N$ equivale a $e^\beta > (N{+}1)/(N{-}1)$,
es decir, $\beta > \ln\!\left((N{+}1)/(N{-}1)\right) = \beta_{\min}$.
$\beta_{\max}$ garantiza la estabilidad local de la dinámica en tiempo discreto:
se obtiene de la condición de que el paso de actualización
$T_s\mu_0 r N \beta/4 < 2$, derivada de la linealización en $\mathbf{x}^*$.
\end{teorema}

\emph{Nota (alcance de las hipótesis).} Al no haber transición de modo, la
coalición no migra al destino y el sistema permanece en la fase estacionaria de
reclutamiento, donde además $\Psi \equiv 1$ (sin descuento espacial). El equilibrio
se analiza, por tanto, únicamente en esa fase; la fase de transporte y el cierre de
la tarea se evalúan empíricamente en los Escenarios~B, C y~E. Las cotas
\eqref{eq:betamin}--\eqref{eq:betamax} son, además, condiciones
\emph{suficientes} derivadas del análisis de Lyapunov: valores de $\beta$ fuera del
intervalo pueden seguir siendo estables, pero el teorema no lo garantiza. El análisis
cubre la fase de reclutamiento homogénea; la navegación del cluster con evitación de
obstáculos durante el transporte queda fuera de la garantía formal y se caracteriza
únicamente de forma empírica, dado que los campos de potencial artificial carecen de
garantías globales frente a mínimos locales.

\begin{proof}[Demostración (esbozo)]
El argumento se construye en cuatro pasos.

\textbf{Paso~1: estructura de juego potencial.}
La función de payoff $f_{ik}(x) = r \cdot \Phi(Nx_k, n)$ es estrictamente
decreciente en $x_k$: cuantos más robots hay en la tarea $k$, menor es el payoff.
Esta propiedad define un juego de congestión simétrico. Por el resultado de
Monderer y Shapley \citep{monderer1996potential}, todo juego de congestión
simétrico es un juego potencial (véase también \citealt{sandholm2010population},
Cap.~3) con función potencial
\[
  V(\mathbf{x}) = \frac{r}{N}\sum_{k=1}^K \int_0^{Nx_k} \Phi(z, n)\,\mathrm{d}z.
\]

\textbf{Paso~2: Smith como ascenso en el potencial.}
Las dinámicas de Smith son un protocolo de revisión de paridad (pairwise protocol)
que en tiempo continuo asciende sobre el potencial: $\dot{V} \geq 0$ a lo largo de
las trayectorias, con igualdad sólo en los equilibrios de Nash
(\citealt{sandholm2010population}, \S5.6 y \S7.1). Para la versión \emph{distribuida}
empleada aquí, \citet{barreiro2017distributed} prueban la estabilidad
asintótica del equilibrio de Nash en juegos potenciales mediante la función de
Lyapunov $E_V(\mathbf{x}) = V(\mathbf{x}^*) - V(\mathbf{x})$, con
$\dot{E}_V = -\mathbf{f}^{\top} L^{(\mathbf{x})}\mathbf{f} \leq 0$. La formulación en
\emph{tiempo discreto} se sustenta en \citet{martinez2020formation,
martinezpiazuelo2022tcss}, que establecen la estabilidad asintótica de las dinámicas
de Smith distribuidas en tiempo discreto bajo una condición suficiente sobre el tamaño
del paso de actualización, inversamente proporcional al número máximo de vecinos en el
grafo.
% TODO-CITA: confirmar nº de teorema vs PDF (Barreiro-Gómez 2017; Martínez-Piazuelo 2020 / 2022a)
Esta cota sobre el paso
es del mismo tipo que la que define $\beta_{\max}$ (Paso~4): la función $\Phi(\cdot,n)$
satisface $|\Phi'(z,n)| \leq \beta/4$, por lo que es Lipschitz-$\tfrac{\beta}{4}$ en
$z$ y el payoff $f_{ik} = r\,\Phi(N\hat{x}_{ik},n)$ es Lipschitz-$\tfrac{rN\beta}{4}$
en $\hat{x}_{ik}$ (el factor $N$ proviene del cambio de variable $z = N\hat{x}_{ik}$),
y es estrictamente decreciente en la masa de la estrategia. Bajo estas hipótesis,
los únicos puntos fijos son los máximos de $V$ sobre el símplex.

\textbf{Paso~3: unicidad del máximo interior.}
La Hessiana de $V$ en el espacio ambiente es diagonal con entradas
$\partial^2 V / \partial x_k^2|_{\mathbf{x}^*} = r N\Phi'(n,n) = -rN\beta/4 < 0$,
donde se usó $\Phi'(z,n) = -\beta\,\Phi\,(1{-}\Phi)$, $\Phi(n,n) = 1/2$, y la
regla de la cadena $\partial (V)/\partial x_k = r\,\Phi(Nx_k,n)$.
En el subespacio tangente al símplex ($\sum_k \mathrm{d}x_k = 0$, dimensión $K{-}1$),
la Hessiana restringida sigue siendo definida negativa pues todos sus valores propios
son $-rN\beta/4$; por tanto, la maximización restringida tiene un único máximo
interior $\mathbf{x}^*$, y en ese punto $Nx_k^* = n$: la restricción de cardinalidad
se satisface exactamente.

\textbf{Paso~4: estabilidad local en tiempo discreto.}
Se define $W(\mathbf{x}) = V(\mathbf{x}^*) - V(\mathbf{x}) \geq 0$ en un entorno
de $\mathbf{x}^*$ (no negativa por ser $\mathbf{x}^*$ el máximo de $V$).
Dado que las dinámicas de Smith ascienden sobre $V$ (Paso~2), $W$ es una
función de Lyapunov estricta en ese entorno.
$\beta_{\min}$ es una \emph{condición de separación de payoff}: garantiza
$\Phi(n{-}1,n)-\Phi(n{+}1,n) > 1/N$, lo que produce pendiente suficiente en $\mathbf{x}^*$
para que $W$ sea localmente estricta.
$\beta_{\max}$ proviene de exigir radio espectral $< 1$ a la linealización de Smith
en $\mathbf{x}^*$: la condición escalar $T_s\mu_0 r N\beta/4 < 2$ es la cota resultante.
Combinados, $\beta \in (\beta_{\min},\beta_{\max})$ asegura que $W$ decrece estrictamente
y $\mathbf{x}^*$ es localmente asintóticamente estable.
\end{proof}

El Teorema~\ref{prop:convergencia} describe un caso idealizado. El sistema
realmente implementado opera con $R$ finito, factor espacial $\Psi$ activo y
transición de modo, condiciones que separan el payoff efectivo del payoff del juego
base. La siguiente proposición acota esa separación.

\begin{proposicion}[Perturbación por consenso local]
\label{prop:perturbacion}
Sea $e_{ik}(t) = \hat{x}_{ik}(t) - x_k(t)$ el error de estimación del DAC, y
supóngase $|e_{ik}(t)| \leq \bar{e}$ para todo $i,k$. Entonces la perturbación del
payoff respecto al juego base está acotada por
\[
  |\delta f_{ik}| = \bigl|f_{ik}(\hat{x}_{ik}) - f_{ik}(x_k)\bigr|
    \;\leq\; \frac{r_k\,\beta\,N}{4}\,\bar{e}.
\]
\end{proposicion}
\begin{proof}
Como $f_{ik} = r_k\,\Phi(N\hat{x}_{ik},n_k)\,\Psi(d_{ik})$ con $\Psi \leq 1$ y
$|\Phi'(z,n)| \leq \beta/4$, el teorema del valor medio aplicado a $\Phi$ con el
cambio de variable $z = N\hat{x}_{ik}$ da
$|\Phi(N\hat{x}_{ik},n_k) - \Phi(Nx_k,n_k)| \leq (\beta/4)\,N\,|e_{ik}|
\leq (\beta N/4)\,\bar{e}$; multiplicando por $r_k\,\Psi \leq r_k$ se obtiene la cota.
\end{proof}

Esta cota acota exclusivamente el canal de error del DAC sobre el payoff. No cubre
formalmente el efecto del factor espacial $\Psi(d)$, el grafo de comunicación móvil
$G(t)$, la transición híbrida reclutamiento--transporte, las saturaciones cinemáticas
ni los mínimos locales del campo de potencial; todos estos efectos se evalúan
empíricamente en los Escenarios~B--E.

En consecuencia, el sistema implementado corresponde a una versión perturbada del
juego base del Teorema~\ref{prop:convergencia}, acotada por la
Proposición~\ref{prop:perturbacion}. La propiedad que cabe esperar es entonces la
\emph{convergencia práctica a una vecindad del equilibrio} (y no la convergencia
exacta al equilibrio Nash ideal), cuyo radio se evalúa empíricamente mediante las
métricas $F$, $D(R)$, $\Delta_{\mathrm{chg}}$ y $\Delta J$.

\begin{observacion}
El Teorema~\ref{prop:convergencia} usa $R = \infty$ y $N = Kn$.
Para los parámetros nominales ($T_s\!=\!0{,}1$, $\mu_0\!=\!0{,}5$, $r\!=\!1$,
$N\!=\!6$, $K\!=\!3$, $n\!=\!2$), los límites~\eqref{eq:betamin}--\eqref{eq:betamax}
dan $\beta_{\min} \approx 0{,}34$ y $\beta_{\max} \approx 26{,}7$, por lo que el
valor nominal $\beta = 2{,}0$ satisface la condición~(C1) con amplio margen.
Cuando $R$ es finito, el error de consenso $e_{ik}$ introduce una perturbación
en los payoffs cuyo efecto sobre el equilibrio se mide experimentalmente en el
Escenario~D mediante la curva de degradación $D(R) = \Theta(R)/\Theta(\infty)$.
La extensión a $N \neq Kn$ (número no divisible exactamente) introduce estados
de inactividad $x_0^* > 0$ y es parte del análisis general que queda como
trabajo futuro. El dominio de atracción de $\mathbf{x}^*$ no se estima en este
esbozo; la convergencia desde inicializaciones arbitrarias sobre el símplex se
verifica numéricamente en el Escenario~A (configuración balanceada).
\end{observacion}

\begin{observacion}[Puente entre fracciones continuas y coaliciones enteras]
\label{obs:fracciones_enteras}
El Teorema~\ref{prop:convergencia} prueba que la \emph{fracción} de robots
converge a $x_k^* = n/N$, un punto del símplex continuo. El requisito de
cardinalidad $|\mathcal{C}_k| \geq n_k$ vive en el dominio \emph{entero}: un robot
se asigna a la tarea $k$ si $s_i = \arg\max_k p_{ik}$. La correspondencia no es
exacta: incluso con $x_k = n/N$, los números enteros $|\mathcal{C}_k|$ pueden
diferir en $\pm 1$ según cómo se resuelvan los empates en $\arg\max$.
En particular, cuando $N = Kn$ y el sistema alcanza el equilibrio interior, la
distribución simétrica implica exactamente $n$ robots por tarea con probabilidad
alta pero no garantizada al~100\,\%. La tasa de factibilidad $F$ (proporción de
episodios en que $|\mathcal{C}_k| \geq n_k$ se alcanza para toda $k$ dentro del
horizonte) es la métrica empírica que cierra este puente.
\end{observacion}

\begin{lema}[Criterio entero de factibilidad]
\label{lema:entero}
Sea $\tilde{x}_k = \tfrac{1}{N}\sum_{i=1}^N \mathbf{1}_{[s_i = k]}$ la ocupación
empírica pura de la tarea $k$, de modo que $N\tilde{x}_k \in \mathbb{Z}_{\geq 0}$ es
el número de robots con estrategia $s_i = k$ (esto es, $N\tilde{x}_k = a_k$). En
general la ocupación física presente cumple
$\mathrm{count}_k(t) = |\mathcal{C}_k(t)| \leq N\tilde{x}_k(t)$, con igualdad cuando todos esos robots se
encuentran dentro del radio de detección $r_{\mathrm{det}}$ de la tarea (coalición ya
formada, régimen estacionario); en transitorios o durante reasignaciones la
desigualdad puede ser estricta. Bajo dicha igualdad:
\textup{(i)}~si $\tilde{x}_k \geq n_k/N$ para toda $k$, entonces
$|\mathcal{C}_k| \geq n_k$ para toda $k$;
\textup{(ii)}~en el caso balanceado $N = Kn$, si $|\tilde{x}_k - n/N| < 1/(2N)$
para toda $k$, entonces $|\mathcal{C}_k| = n$ exactamente.
\end{lema}
\begin{proof}
Bajo la igualdad $|\mathcal{C}_k| = N\tilde{x}_k$ y siendo $N\tilde{x}_k$ entero:
(i) $\tilde{x}_k \geq n_k/N$ implica $N\tilde{x}_k \geq n_k$; (ii) $|N\tilde{x}_k - n| < 1/2$
con $N\tilde{x}_k$ y $n$ enteros fuerza $N\tilde{x}_k = n$.
\end{proof}
Este lema conecta el equilibrio continuo $x_k^* = n/N$ del
Teorema~\ref{prop:convergencia} con la factibilidad entera: en régimen estacionario,
una vecindad de radio $1/(2N)$ alrededor de $x_k^*$ es suficiente para la cardinalidad
exacta.

\begin{observacion}[Tiempo de convergencia del DAC]
Para los parámetros nominales, el tiempo de contracción del error de seguimiento
del DAC se estima a partir del término de difusión en la red: el error
$\hat{x}(t{+}1) - \bar{x}(t{+}1) = (I - \varepsilon L)(\hat{x}(t) - \bar{x}(t))$
decae como $(1 - \varepsilon\lambda_2)^t$, donde $\varepsilon = 0{,}05$ y $\lambda_2$
es el valor de Fiedler del grafo de comunicación, que en el Escenario~A depende de
las posiciones 2D y del rango $R$. Tomando un valor representativo $\lambda_2 \approx 1$
(orden de magnitud del Fiedler en configuraciones densas en 2D), el número de pasos
para reducir el error residual en un factor $1/e$ es:
\[
  \tau_{\mathrm{pasos}} \approx \frac{1}{\varepsilon \cdot \lambda_2}
  = \frac{1}{0{,}05 \cdot 1} = 20\,\text{pasos},
  \quad
  \tau_{\mathrm{conv}} \approx \tau_{\mathrm{pasos}} \cdot T_s = 2\,\text{s}.
\]
En el caso límite de comunicación global ($R = \infty$) el grafo es completo y
$\lambda_2 = N = 6$, lo que reduce $\tau_{\mathrm{pasos}}$ a ${\approx}\,3$ pasos.
Este tiempo (${\approx}20\,T_s$) es del mismo orden que el tiempo de formación de
coalición esperado, lo que plantea una tensión de diseño: el lazo de Smith
opera con estimaciones que aún no han alcanzado la media global exacta.
El término $\Delta y_{ik}$ del DAC compensa el salto local instantáneamente,
pero el error de seguimiento acotado por $O(K/N)$ (véase Obs.~\ref{obs:init_obs})
persiste $\sim 20$ pasos adicionales hasta que la información se propaga por la red.
Durante ese transitorio, el payoff $f_{ik}$ puede sobrestimar la ocupación (sesgo
estabilizador en saturación) o subestimarla (perturbador en reclutamiento).
El efecto neto se cuantifica mediante $\Delta_{\mathrm{chg}}$ en el Escenario~A;
la ablación de $\hat{x}$ ideal frente a estimado aísla el impacto del error
de consenso sobre el equilibrio resultante.
\end{observacion}

\subsection{Diseño experimental}
\label{subsec:experimental}

El diseño compara cinco tratamientos en seis escenarios. Todos los tratamientos
usan los mismos valores de parámetros físicos, las mismas condiciones iniciales
de posición (controladas por semilla) y la misma secuencia de aparición de tareas.

\subsubsection{Tratamientos}

Los tratamientos forman una escalera incremental: cada uno añade exactamente un
mecanismo respecto al anterior, de modo que la diferencia de desempeño entre dos
niveles consecutivos es atribuible a ese mecanismo. En particular, T3 y T4
comparten todo lo demás (el mismo payoff \eqref{eq:payoff}, el mismo estimador
DAC \eqref{eq:consenso} para $\hat{x}_{ik}$ y los mismos parámetros
$\beta,\lambda,\mu_0$) y difieren únicamente en la ley de revisión, de manera
que la comparación T3 frente a T4 aísla el efecto de la política de revisión.

\begin{itemize}
  \item \textbf{T1 -- Heurística voraz local}: cada robot selecciona la tarea más
    cercana que todavía necesita robots ($|\mathcal{C}_k| < n_k$), usando solo
    información geométrica local. No hay estimación de ocupación ni dinámica de
    revisión: mide si la sola cercanía geométrica resuelve la asignación.

  \item \textbf{T2 -- DAC con regla voraz}: se ejecuta el consenso dinámico de
    media \eqref{eq:consenso} para estimar la ocupación $\hat{x}_{ik}$, pero la
    asignación aplica una regla voraz sobre esas estimaciones, sin dinámica
    poblacional de revisión. Aísla si basta con estimar bien la ocupación, sin una
    dinámica de decisión que la explote.

  \item \textbf{T3 -- Replicador distribuido modificado}: dinámica del replicador
    en tiempo discreto \citep{quijano2017population} con el mismo payoff y el
    mismo estimador DAC que T4. La actualización
    \[
      p_{ik}(t{+}1) = p_{ik}(t) + T_s\,\mu_i(t)\,p_{ik}(t)\bigl(f_{ik}(t) - \bar{f}_i(t)\bigr),
    \]
    proyectada al símplex, requiere la media de payoff de la población
    $\bar{f}_i$, que no es una cantidad local: se estima mediante una
    \emph{segunda capa de consenso} sobre los valores $\bar{f}_j$ de los vecinos
    (véase \S\,\ref{subsec:porquesmith}). Esta arquitectura de doble consenso es
    necesaria para que la comparación con T4 sea justa: calcular $\bar{f}_i$ de
    forma puramente local eliminaría precisamente el sesgo topológico que T3 está
    diseñado para exponer. Examina si una dinámica poblacional basada en la media de
    payoff resuelve el problema.

  \item \textbf{T4 -- Dinámicas de Smith con revisión modulada (método propuesto)}:
    tres dinámicas acopladas con función de payoff \eqref{eq:payoff}, dinámica de
    Smith \eqref{eq:smith_update}, cuya comparación bilateral evita estimar la
    media global, y tasa de revisión espacialmente modulada
    \eqref{eq:revision_rate}. El contraste entre $\mu_i$ modulada y $\mu_i=\mu_0$
    constante constituye la ablación interna del método. Evalúa si la comparación
    bilateral de Smith mejora sobre el replicador al reducir $R$.

  \item \textbf{T5 -- Referencia centralizada replanificada}: con acceso al
    estado global, se resuelve la asignación de mínima distancia respetando la
    cardinalidad: cada tarea $k$ se expande en $n_k$ \emph{slots} idénticos; se
    añaden \emph{filas ficticias} si $N < \sum_k n_k$ o \emph{columnas ficticias}
    si $N > \sum_k n_k$ (al menos un robot queda inactivo); y se minimiza la
    distancia euclídea total robot--slot mediante el algoritmo húngaro sobre la
    matriz cuadrada resultante. Emplea la misma cinemática uniciclo y los mismos
    límites de velocidad ($v_{\max},\omega_{\max}$) que el resto de tratamientos.
    La reasignación se replanifica cada $\Delta T_c = 1\,\text{s}$ y ante eventos
    (aparición o desaparición de tareas, retirada de robots). La función objetivo,
    minimizar la distancia de asignación, no coincide necesariamente con maximizar
    $\Theta$, por lo que T5 es una referencia de eficiencia de asignación, no un
    límite superior absoluto del throughput. Cuantifica cuánto se pierde por
    descentralizar.
\end{itemize}

\subsubsection{Escenarios}

La Tabla~\ref{tab:escenarios} describe los seis escenarios del núcleo experimental
(A--F) y un escenario exploratorio adicional (G, \S\,\ref{subsec:heterogenea}). Los
seis del núcleo se ejecutan con $m = 30$ repeticiones por tratamiento con semillas
distintas pero fijadas.

\begin{table}[htbp]
\centering
\small
\renewcommand{\arraystretch}{1.35}
\caption[Escenarios experimentales]%
{Escenarios experimentales: configuración, propósito y parámetros clave.}
\label{tab:escenarios}
\begin{tabularx}{\textwidth}{@{}>{\bfseries\sffamily\color{VIUOrange}}p{0.07\textwidth} p{0.30\textwidth} Y@{}}
\toprule
\normalfont\bfseries\color{VIUDark}Esc. &
  \bfseries Configuración &
  \bfseries Propósito \\
\midrule
A &
  6 AGV, 3 tareas. \textbf{(a)}~$n_k = 2$ (balanceado, $N=Kn=6$);
  \textbf{(b)}~$n_k = 1$ (línea base). &
  (a)~Validación numérica de H1 con $\Psi\equiv 1$, $R=\infty$, $N=Kn$.
  (b)~Línea base sin requisito de coalición. \\
\addlinespace[2pt]
B &
  6 AGV; 1 tarea ligera ($n=1$) + 1 tarea pesada ($n=3$). &
  Prueba central de coalición mínima: 1 robot a la ligera, 3 a la pesada. \\
\addlinespace[2pt]
C &
  12 AGV, 5 tareas con $n_k \in \{1,1,2,3,4\}$. Recompensas
  $r_k \in \{2{,}0,\, 1{,}0,\, 0{,}5\}$. &
  Competencia y prioridad: tareas de $r_k$ alto atraen agentes antes,
  sin coordinación central. \\
\addlinespace[2pt]
D &
  15 AGV, 5 tareas ($n_k \in \{1,2,2,3,4\}$). $R$ barrido en 10 puntos
  log-espaciados de 1\,m a 15\,m, más $R=\infty$. &
  Curva de degradación $D(R) = \Theta(R)/\Theta(\infty)$, concentrada
  en la zona de transición de conectividad. \\
\addlinespace[2pt]
E &
  10 AGV, 3 tareas ($n_k \in \{2,3,4\}$). Eliminación abrupta de 1 robot
  en $t=20\,\text{s}$ y 2 tareas nuevas ($n_k \in \{1,2\}$) en $t=30\,\text{s}$. &
  Reactividad ante perturbación: tiempo de reconstitución de coalición tras
  cambios abruptos en el conjunto. \\
\addlinespace[2pt]
F &
  4--6 Pioneer~3-DX en CoppeliaSim. Escenario B o C con carga pasiva
  acoplada cinemáticamente al centroide. &
  Validación cualitativa de la formación de coalición y la navegación
  cooperativa. \\
\addlinespace[2pt]
G &
  \emph{Exploratorio.} Flota con capacidades heterogéneas (p.\,ej.\ $c_i \in \{30,30,50\}$\,kg)
  y cargas con demandas de capacidad mixtas (p.\,ej.\ $M_k \in \{10,100\}$\,kg). &
  Exploración cualitativa de la extensión de capacidad heterogénea
  (\S\,\ref{subsec:heterogenea}); de carácter exploratorio, sus resultados no forman
  parte de los criterios de aceptación. \\
\bottomrule
\end{tabularx}
\end{table}

\begin{table}[htbp]
\centering
\small
\caption[Hipótesis, comparaciones, escenarios y métricas]{Correspondencia entre
  hipótesis, comparación principal, escenarios de contraste y métrica principal.}
\label{tab:hipescenario}
\begin{tabularx}{\textwidth}{>{\bfseries}p{0.07\textwidth}p{0.30\textwidth}p{0.10\textwidth}Y}
\toprule
Hip. & Comparación principal & Escenarios & Métrica principal \\
\midrule
H1 & T2 vs.\ T3/T4 & B, C & $F$, $T_{\mathrm{coal}}$, $\Delta_{\mathrm{chg}}$ \\
H2 & T2 vs.\ T3/T4 & B, C & $F$, $\Theta$ \\
H3 & T3 con barrido de $R$ & D & $F$, $\Theta$, $D(R)$ \\
H4 & T3 vs.\ T4 & C, D & $F$, $\Theta$, $\Delta_{\mathrm{chg}}$ \\
H5 & T4($\mu$ constante) vs.\ T4($\mu$ modulada), ablación & A, B &
  $\Delta_{\mathrm{chg}}$, $T_{\mathrm{coal}}$ \\
H6 & T4 vs.\ T5 & C, D, E & $\Delta J$, $\Theta$ \\
\bottomrule
\end{tabularx}
\end{table}

\subsection{Métricas}

Todas las métricas se definen antes de ejecutar la campaña experimental y se
calculan automáticamente sobre los registros de cada episodio.

\begin{table}[htbp]
\centering
\caption[Métricas de evaluación]{Métricas de evaluación: definición e interpretación.}
\label{tab:metricas}
\begin{tabularx}{\textwidth}{>{\bfseries}p{0.08\textwidth}p{0.24\textwidth}Y}
\toprule
Símbolo & Nombre & Definición / Interpretación \\
\midrule
$\Theta$ & Tasa de throughput &
  Tareas completadas por minuto promediadas sobre el episodio. \\
$\bar{T}$ & Tiempo medio de entrega &
  Promedio del tiempo entre aparición de la carga y entrega al destino. \\
$T_{\text{coal}}$ & Tiempo de formación de coalición &
  Tiempo entre aparición de la tarea y el instante en que
  $|\mathcal{C}_k(t)| \geq n_k$ por primera vez. \\
$C$ & Sobrecarga de comunicación &
  Mensajes intercambiados por robot por paso temporal. \\
$D(R)$ & Ratio de degradación &
  $\Theta(R)/\Theta(\infty)$; mide pérdida relativa de rendimiento con $R$ finito. \\
$F$ & Tasa de factibilidad &
  Porcentaje de tareas que alcanzan $|\mathcal{C}_k| \geq n_k$ dentro del
  horizonte temporal $T$. \\
$\Delta_{\text{chg}}$ & Cambios de asignación &
  Cambios de estrategia por robot por paso temporal:
  $\frac{1}{T}\sum_{t=1}^{T}\mathbf{1}_{[s_i(t)\neq s_i(t-1)]}$;
  la normalización por $T$ hace comparables episodios de distinta duración. \\
$D_{\text{tot}}$ & Distancia total recorrida &
  $\sum_i \int_0^T \norm{v_i(t)}\,\mathrm{d}t$; proxy de consumo energético
  de la flota. \\
$\Delta J$ & Brecha de descentralización &
  $(\Theta_{\text{cent}} - \Theta_{\text{dist}}) / \max(\Theta_{\text{cent}}, \delta_0)$,
  con $J \triangleq \Theta$ y $\delta_0 = 10^{-3}$;
  mide cuánto pierde el distribuido respecto al centralizado ($\Delta J = 0$: sin pérdida). \\
$E_{\text{DAC}}$ & Error medio de estimación &
  $\frac{1}{NKT}\sum_{t}\sum_{i}\sum_{k}|\hat{x}_{ik}(t) - x_k(t)|$;
  aísla empíricamente el efecto del error de consenso del efecto de la dinámica de
  revisión. \\
$N_{\text{fail}}$ & Episodios fallidos &
  Número de episodios por escenario en que alguna tarea no alcanza
  $|\mathcal{C}_k| \geq n_k$ dentro del horizonte; se expresa como tasa de fallo
  por tratamiento. \\
\bottomrule
\end{tabularx}
\end{table}

Los resultados se presentan como media $\pm$ desviación estándar sobre $m = 30$
repeticiones. Las comparaciones pareadas entre tratamientos (p.ej., dinámica de
Smith frente a líneas base, $\mu_i$ modulada frente a constante) se contrastarán
mediante pruebas no paramétricas: la prueba de Wilcoxon de rangos con signo (datos
pareados) o la U de Mann-Whitney (dos muestras independientes). El análisis estadístico
de la entrega final contemplará, además, corrección para comparaciones múltiples,
intervalos de confianza por bootstrap y medidas de tamaño de efecto; el detalle se fija
en la fase de resultados. Se incluye un episodio fallido por
escenario para documentar los límites del método. El registro completo de
trayectorias $\mathbf{p}_i(t)$ para todo $i$ y todo $t$ se almacena para
análisis a posteriori de la dinámica de preferencias.

\subsection{Implementación en Python}

La plataforma de simulación está implementada en Python con una arquitectura
modular que separa el modelo del dominio, las dinámicas distribuidas, los
tratamientos de comparación, el motor de simulación y las métricas. Las
operaciones numéricas críticas (consenso, proyección al símplex, cálculo del
gradiente) se vectorizan con NumPy. La Tabla~\ref{tab:params} recoge los
parámetros del sistema junto con sus valores nominales y el rango de variación
considerado en los experimentos.

\begin{table}[htbp]
\centering
\caption[Parámetros del sistema]{Parámetros del sistema y valores nominales para los experimentos.}
\label{tab:params}
\begin{tabularx}{\textwidth}{>{\ttfamily}p{0.12\textwidth}p{0.28\textwidth}p{0.10\textwidth}Y}
\toprule
\normalfont Símbolo & Parámetro & Valor nominal & Rango / Nota \\
\midrule
$T_s$ & Período de muestreo & 0{,}10 s & Fijo \\
$\beta$ & Pendiente sigmoide & 2{,}0 & (1{,}0, 5{,}0) \\
$\lambda$ & Escala espacial del payoff & 5{,}0 m & (2{,}0, 10{,}0) m \\
$r_0$ & Payoff de inactividad & 0{,}05 & Fijo \\
$\mu_0$ & Tasa de revisión base & 0{,}50 & (0{,}1, 1{,}0) \\
$\mu_{\min}$ & Piso de revisión ($= 0{,}02\,\mu_0$) & 0{,}01 & Fijo \\
$\sigma_{\text{rev}}$ & Ancho espacial de revisión & 2{,}0 m & (1{,}0, 5{,}0) m \\
$\tau_s$ & Umbral de modo transporte & 3 pasos & Fijo \\
$\rho_s$ & Umbral de histéresis para cambio de estrategia & 0{,}05 & Fijo \\
$\varepsilon$ & Ganancia de consenso & 0{,}05 & $(0, 1/d_{\max})$ \\
$R$ & Rango de comunicación & $\infty$ & (1{,}0 m, $\infty$) \\
$\alpha$ & Ganancia de atracción & 1{,}0 & Fijo \\
$\eta$ & Ganancia de repulsión & 0{,}5 & Fijo \\
$\delta$ & Ganancia de formación & 0{,}8 & Fijo \\
$r_{\mathrm{det}}$ & Radio de detección de coalición & 1{,}0 m & Fijo \\
$r_{\mathrm{form}}$ & Radio de activación de formación & 1{,}5 m & Fijo \\
$r_{\mathrm{robot}}$ & Radio físico del AGV & 0{,}25 m & Pioneer~3-DX \\
$r_{\mathrm{load}}$ & Radio de la carga & 0{,}30 m & Valor nominal \\
$\varepsilon_r$ & Regularización de repulsión & 0{,}01 m & Fijo \\
$\varepsilon_f$ & Regularización de formación & 0{,}01 m & Fijo \\
$\eta_o$ & Ganancia de repulsión de obstáculos & 0{,}5 & valor nominal, ajustable \\
$\varepsilon_o$ & Regularización de obstáculos & 0{,}05 m & valor nominal, ajustable \\
$\omega_{\max}$ & Velocidad angular máxima & 2{,}5 rad/s & Pioneer~3-DX \\
$v_{\max}$ & Velocidad máxima & 0{,}5 m/s & Límite conservador por seguridad en escena reducida; el Pioneer~3-DX admite ${\sim}1{,}2\,\text{m/s}$ \\
$\ell$ & Distancia look-ahead & 0{,}15 m & Pioneer~3-DX \\
\bottomrule
\end{tabularx}
\end{table}

Cada experimento se parametriza mediante un descriptor declarativo que fija
todos los parámetros y las semillas aleatorias, de modo que los resultados son
reproducibles. Durante la simulación se registran las trayectorias
$\mathbf{p}_i(t)$, $\mathbf{q}_i(t)$ y $\hat{\mathbf{x}}_i(t)$ en cada paso
temporal; las métricas se calculan en una etapa de postprocesamiento
posterior.

El código, las configuraciones YAML, las semillas y los registros de trayectorias,
asignaciones, grafos, payoffs y métricas se organizan en un repositorio versionado;
para una publicación posterior se prevé una versión archivada con DOI.

\subsection{Validación en CoppeliaSim}

El Escenario~F replica uno de los escenarios de coalición mínima (B o C) con
robots Pioneer~3-DX en CoppeliaSim. Esta validación no añade métricas cuantitativas;
confirma que el comportamiento colectivo observado en Python es coherente con la
física del simulador robótico: los robots se desplazan hacia las tareas correctas,
forman la geometría de coalición y reaccionan ante una perturbación (retirada de un
robot a mitad de la misión). La escena incluye un almacén simplificado con zonas de
carga marcadas en el suelo, compatible con el tamaño operativo de los Pioneer~3-DX.

El alcance de esta validación es acotado: la simulación en CoppeliaSim verifica la
viabilidad cinemática del esquema de asignación, formación y navegación con robots
Pioneer~3-DX (saturación de velocidades, geometría espacial). No abarca la
manipulación cooperativa por contacto ni la distribución de fuerzas sobre la carga,
que pertenecen a una capa de control complementaria
\citep{ebel2024cooperative, rosenfelder2024force} fuera del alcance de este TFM. La
carga se representa como un cuboide pasivo acoplado cinemáticamente al centroide de
la coalición una vez formada.

\subsection{Criterios de aceptación del método}

El método se considerará satisfactorio si:

\begin{enumerate}
  \item La tasa de factibilidad del tratamiento~T4 (método propuesto) es
    $F \geq 0{,}85$ en los Escenarios~B y~C.
  \item El tiempo de formación de coalición $T_{\text{coal}}$ del tratamiento~T4
    es menor que el del tratamiento~T1 en al menos el 80\,\% de las repeticiones
    del Escenario~B.
  \item El número de cambios de asignación $\Delta_{\text{chg}}$ del
    tratamiento~T4 con tasa espacial \eqref{eq:revision_rate} es inferior al del
    mismo tratamiento con tasa constante $\mu_i = \mu_0$, validando la
    Hipótesis~5 (ablación de la tasa de revisión).
  \item La validación numérica del Teorema~\ref{prop:convergencia}
    en el Escenario~A (configuración balanceada $n_k=2$, $N=Kn=6$) con 30 episodios
    muestra convergencia ($T_{\mathrm{conv}} < \infty$ y
    $|x_k(T_{\text{conv}}) - x_k^*| < 0{,}01$ para todo $k$) en al menos 27 episodios
    (tasa de éxito ${\geq}90\%$).
  \item La escena de CoppeliaSim (Escenario~F) muestra visualmente la formación
    de la coalición requerida y la reacción ante perturbación.
\end{enumerate}

Los criterios 1--5 son \textbf{conjuntivos}: deben cumplirse todos para declarar el
método satisfactorio. El criterio~5 es condición necesaria pero de carácter
cualitativo; los criterios~1--4 son cuantitativos y verificables estadísticamente.
Si alguno de los criterios~1--4 no se cumple, el experimento se replicará
ajustando los parámetros nominales dentro del rango indicado en la
Tabla~\ref{tab:params} antes de concluir el rechazo. El Escenario~G es exploratorio y
queda excluido de los criterios de aceptación conjuntivos anteriores.

\subsection{Extensión a capacidad heterogénea (formulación)}
\label{subsec:heterogenea}

Se considera la generalización en que cada robot $i$ posee una capacidad $c_i > 0$ y
cada carga $k$ exige una capacidad acumulada $M_k$. Una coalición $\mathcal{C}_k$ es
factible si $\sum_{i\in\mathcal{C}_k} c_i \geq M_k$. El payoff de cardinalidad
$\Phi(N\hat{x}_{ik}, n_k)$ se reemplaza por una demanda de capacidad
$\Phi_c(\hat{c}_k, M_k)$, monótonamente decreciente en la capacidad estimada ya
comprometida $\hat{c}_k$, estimable por consenso de forma análoga a $\hat{x}_{ik}$.

Esta extensión rompe la simetría explotada en el Teorema~\ref{prop:convergencia}: con
capacidades heterogéneas el juego deja de ser simétrico, el equilibrio Nash deja de
ser único e interior, y el conjunto de coaliciones factibles
$\{\mathcal{C}_k : \sum_{i\in\mathcal{C}_k} c_i \geq M_k\}$ es un problema combinatorio
de tipo \emph{subset-sum}, NP-duro en general. Resolverlo de forma distribuida mediante
dinámicas poblacionales (incluido el riesgo de coaliciones infactibles localmente
estables) es un problema abierto. Por ello se presenta solo a nivel
de formulación; su análisis formal de convergencia y su validación exhaustiva quedan
fuera del alcance de este TFM (véase \S\,\ref{sec:conclusiones}, trabajo futuro). El
Escenario~G explora cualitativamente su comportamiento.


%% ====================================================================
\section{Marco teórico}
\label{sec:marco}
%% ====================================================================

\textit{[Sección reservada para la entrega final. Se desarrollarán los siguientes
apartados:]}

\begin{itemize}
  \item 7.1 Juegos poblacionales y equilibrio Nash
  \item 7.2 Dinámicas de Smith: definición, propiedades y comparación con
    el replicador
  \item 7.3 Juegos de congestión ponderados: estructura potencial
  \item 7.4 Consenso distribuido en grafos dinámicos
  \item 7.5 Transporte cooperativo multi-robot: estado del arte
  \item 7.6 Formación de coaliciones en sistemas multi-agente
\end{itemize}


%% ====================================================================
\section{Resultados y análisis}
\label{sec:resultados}
%% ====================================================================

\textit{[Sección reservada para la entrega final. Los resultados se generarán
desde la plataforma Python mediante la campaña experimental descrita en la
Sección~\ref{subsec:experimental}. Se incluirán:]}

\begin{itemize}
  \item Validación numérica de la Proposición~\ref{prop:convergencia}
    (convergencia al equilibrio en el Escenario~A, configuración balanceada $n_k=2$).
  \item Tablas comparativas de todas las métricas para los seis escenarios y
    cinco tratamientos.
  \item Curva de degradación $D(R)$ del Escenario~D.
  \item Trayectorias de preferencias $\mathbf{p}_i(t)$ mostrando la formación
    dinámica de la coalición.
  \item Imágenes y secuencias de CoppeliaSim del Escenario~F.
\end{itemize}


%% ====================================================================
\section{Conclusiones y recomendaciones}
\label{sec:conclusiones}
%% ====================================================================

\textit{[Sección reservada para la entrega final. Las conclusiones discutirán:]}

\begin{itemize}
  \item Cumplimiento de los criterios de aceptación y contraste de hipótesis.
  \item Condiciones bajo las cuales el método propuesto mejora o no a los
    tratamientos de referencia.
  \item Limitaciones observadas en los experimentos: efecto del grafo irregular,
    escenarios con $N \neq Kn$ y robustez ante perturbaciones.
  \item Líneas de trabajo futuro: extensión doctoral DRD-PI, tasas de revisión
    dependientes del grafo, generación de tareas en tiempo real.
  \item Extensión a capacidad heterogénea: reemplazo de la cardinalidad mínima por la
    restricción de capacidad acumulada $\sum c_i \geq M_k$, pérdida de la simetría del
    juego, conexión con \emph{subset-sum} distribuido y juegos poblacionales multi-clase.
\end{itemize}


%% ====================================================================
\bibliographystyle{apalike}
\bibliography{references}
%% ====================================================================

\end{document}
